% CHEMSA LIVING MANUSCRIPT
% WORKING / NOT A RELEASE / NOT SUBMISSION-READY
% Seeded from the user-supplied current source: ChemSA_Combined_Main (9).tex
% Program governance: SymC General Operations Manual v0.8.0
% Scientific edits must remain synchronized with MANUSCRIPT_STATUS.json and repo evidence.

\documentclass[11pt]{article}
\usepackage[margin=1in]{geometry}
\usepackage{natbib}
\usepackage{booktabs}
\usepackage[hidelinks,pdfusetitle]{hyperref}
\usepackage{tabularx}
\usepackage{amsmath,amssymb}
\usepackage{graphicx}
\usepackage{microtype}

\renewcommand{\arraystretch}{1.2}

\title{Generator-First Classification of Damped Vibrational Modes,\\
and a Three-Step Procedure for Interpreting Spectral Linewidths as Mechanical Damping}

\author{Nate Christensen\\ SymC Universe Project, Missouri, USA\\ nate.christensen@symcuniverse.org}
\date{}

\begin{document}
\maketitle

\begin{abstract}
A repeated vibrational eigenvalue and a measured spectral linewidth are both observables that do not describe themselves, and ChemSA addresses one of each. A repeated eigenvalue may be semisimple, with independent eigenvectors, or defective, forming an exceptional point; the two are identical in the eigenvalues and differ only in the eigenvectors. ChemSA is a generator-first classifier that distinguishes them by testing eigenvector deficiency, scopes every exceptional-point verdict to the declared provenance of the supplied equation, and returns crowded unresolvable clusters as unresolved. A constructive counterexample (identical local architecture, activated rates differing by a factor of about 148 under a common unspecified prefactor) keeps local mode architecture from being read as a rate predictor.

The same discipline is then applied to linewidth interpretation. A measured spectral line may combine static inhomogeneous broadening with a homogeneous width, and the homogeneous width may itself contain lifetime relaxation, pure dephasing, and orientational relaxation; only the lifetime contribution is a mechanical damping coefficient. This decomposition is standard and is not introduced here; a 2025 survey reports that incorrect width-to-lifetime conversions remain prevalent, producing errors of factors of two or \(\pi\), with pure dephasing and inhomogeneous broadening frequently omitted \citep{dubajic2025}. A three-step procedure separates these: isolate a mode-specific homogeneous response; separate its lifetime, pure-dephasing, and orientational contributions using an independent lifetime measurement and an orientational term that is measured, bounded, or independently excluded; and only then compute $\chi = (1/T_1)/2\omega_0$ from the lifetime. The ratio $R_\phi = (1/T_\phi)/(1/2T_1)$ quantifies the bias ($1 + R_\phi + R_{\mathrm{or}}$) of reading the total homogeneous width as a lifetime proxy. In the Rh(CO)$_2$acac example, the source independently reports echo and pump-probe measurements and attributes the strong temperature-dependent homogeneous broadening primarily to pure dephasing; the homogeneous linewidth grows by a factor of about 13.6 from 3.4 to 250~K, which is the one quantitative result here, being derived from two source-text-anchored widths. The lifetime-derived $\chi \approx 3\times10^{-5}$ shown alongside it is illustrative, pending calibrated pixel extraction of the plotted $T_1$ values, since every pump-probe reading used here is a by-eye reading of the source figure.
\end{abstract}

\vspace{0.5em}

\section{Introduction}

\subsection{A degeneracy is not self-describing}

When two modes of a chemical system approach the same frequency and linewidth, a spectrum alone cannot say what kind of degeneracy it is. Two possibilities look identical in the eigenvalues but differ completely in the physics. In a semisimple (diagonalizable) degeneracy the two modes share an eigenvalue but keep independent eigenvectors: perturb the system and they separate smoothly. In an exceptional point the eigenvectors themselves coalesce; the system becomes defective, and it responds to perturbation with the square-root sensitivity that makes exceptional points valuable for sensing and dangerous to misidentify. The two are distinguished not by how close the eigenvalues are but by whether the eigenvectors have collapsed, a property a frequency-and-linewidth fit does not report.

A second trap is layered on the first. A degeneracy read from a reduced description, an open-system or effective equation where the environment has been integrated out, is not the same object as one in the full generator. A repeated root in a Lindblad-reduced equation can look like an exceptional point while corresponding to nothing defective in the underlying closed system. Certifying a physical exceptional point therefore requires knowing the provenance of the equation it came from, and many reported degeneracies do not carry that provenance.

\subsection{The machinery first}

The primary contribution is a classification engine that closes both traps, together with a scope fence that keeps it honest. ChemSA distinguishes a defective degeneracy from a semisimple one by testing eigenvector deficiency directly, returning different verdicts for two systems with the same repeated eigenvalue (Section~2.3); it gates the exceptional-point verdict on declared provenance, so defectiveness found in a reduced equation is reported as an exceptional point of that supplied representation with the parent status left explicitly unresolved; and it is bounded by a constructive counterexample, two systems with identical local architecture and activated rates differing by a factor of about 148 under a common unspecified prefactor (Section~3). The current ChemSA interface contains no reaction-rate estimator. The constructive counterexample of Section~3 is the mathematical statement of why local spectral architecture alone cannot supply one; the regression test is a software guard that keeps the implementation from drifting back across that boundary. The proof bounds the claim, the test bounds the code. These three core scientific safeguards, together with three implementation validity gates (scalar-\(\chi\) applicability, source gating of numerical inputs, and finite-precision equality caution; Section~2.3) are the defensible core: for a single published damped-oscillator fit the damping ratio is already known, so a catalogue of such ratios adds little; what the engine adds is the exceptional-point-versus-semisimple discrimination, the provenance gate, and the scope fence, none of which a standard fit provides.

\subsection{The coordinate, and the boundary of where it is a damping ratio}

To place a mode's stability on an axis the engine uses the damping ratio $\chi$ = \(\gamma/2\omega_0\), read from the poles of the linearized generator. What $\chi$ is must be stated carefully, because it is easy to overreach.

The mechanical damping ratio is defined for a second-order modal representation: a scalar \(M,C,K\) mode, or a system whose damping is proportional so that it decouples into independent modes (the Caughey--O'Kelly condition \citep{caughey}). That condition is TESTED by the engine, not accepted on declaration: \texttt{scalar\_reduction\_valid} computes the relative mass-metric commutator of \(M^{-1}C\) and \(M^{-1}K\) and withholds the mechanical \(\chi\), with a stated reason, whenever it is nonzero. Where the test passes, the modal basis is fixed by diagonalizing the damping within each degenerate stiffness block, since a repeated stiffness eigenvalue leaves the basis undetermined and only one rotation inside that subspace decouples the damping as well. In that representation \(\chi = \gamma/2\omega_0\) is a genuine mechanical damping ratio: $\chi$ \(<1\) rings, $\chi$ \(=1\) is critical, $\chi$ \(>1\) is two real decay timescales. This is the regime in which $\chi$ is meaningful as a continuous mechanical coordinate.

For a general first-order generator (an arbitrary kinetic Jacobian, a non-modal coupled system) no mechanical $\chi$ is claimed. What the engine reports there is a spectral pole-angle index \(\rho = -\mathrm{Re}(\lambda)/|\lambda|\) together with a categorical temporal class (oscillatory, repeated-root, or monotone). The distinction matters and a careful reader will insist on it: \(\rho\) is the cosine of the pole's angle, and nothing more. It EQUALS a mechanical damping ratio only when the pole belongs to a scalar or classically damped second-order factor that has been separately licensed; for a general first-order generator there is no mechanical \(\chi\) for it to equal, and reading \(\rho\) as one is the precise conflation this work exists to prevent. For a stable real pole \(\rho = 1\) identically, for every stable monotone system, no matter how far from any critical point, so \(\rho\) cannot by itself distinguish a repeated-root mode from ordinary stable real relaxation; only the second-order modal picture recovers that distinction. Because that value carries no information, ChemSA does not compute or report it: \(\rho\) is returned only for a complex pole, and is omitted (not set to 1) for a real one, so its presence in a record already signals which case applies. Accordingly, ChemSA assigns a true $\chi$ only when a second-order or decoupled modal representation is established, and otherwise reports pole geometry and temporal class without forcing a mechanical $\chi$.

What linearization does give, stated honestly, is a common spectral footing: disparate systems reduce to poles, and the engine reads those poles the same way. That is what lets one demonstration span phases, a property of linearization, not a claim that one number governs disparate physics, and not a claim that arbitrary first-order networks occupy the same continuous under-, critical-, over-damped axis.

\subsection{What this paper does not claim}

This work makes no claim about reaction rate, yield, commitment probability, or selectivity: the current interface exposes no rate estimator, and the scope proof of Section~3 shows why local classification cannot determine rate (identical local architecture, rates differing by 148). It does not claim a real system that rests at $\chi$ \(=1\): the critical point is a boundary systems cross rather than a state a real system occupies. The linewidth hierarchy of Part~III does not claim a coupled-mode mechanical catalogue: its mechanical-placement claim is confined to a single reconstructed trajectory in which an independent lifetime measurement is available, and a chemical system whose coupled generator sits on an exceptional point remains open work, awaiting two linewidths that only a dedicated experiment can supply.

\subsection{Relation to prior work}

The scalar damping coordinate used here is not introduced by this paper. Earlier work established \(\chi = \gamma/2\omega_0\) as the structural boundary of a dominant second-order response, showed that \(\chi = 1\) marks a boundary that systems cross rather than a state they occupy or optimize toward, restricted exceptional-point terminology to defective realizations, and noted that a strongly multimode system cannot be compressed into a single scalar \(\chi\) \citep{stabilityarch}. None of those results is re-claimed here. What the present work adds is that each of them becomes a test that can refuse: the restriction to a scalar response becomes a Caughey--O'Kelly check that withholds \(\chi\) with a stated reason, the restriction of exceptional-point terminology to defective realizations becomes an algebraic-versus-geometric multiplicity test, and the caution that a reduced description does not establish its parent becomes the provenance gate. Section~3 adds one result in the other direction: the boundary is specifically a property of a stable well, and no counterpart exists at an inverted mode.

What the present work adds is the machinery that those statements require once the system is not scalar. In a scalar second-order mode, \(\chi = 1\) produces a defective companion block, so the critical point and the exceptional point coincide and neither has to be tested for separately. In a coupled spectrum that identity fails: an eigenvalue coincidence may be defective or semisimple, and the two are indistinguishable in the eigenvalues alone, so defectiveness becomes something to test rather than something to infer. This paper therefore generalizes the treatment from a dominant scalar mode to coupled quadratic pencils and general first-order generators, and supplies what that generalization needs: an executable algebraic-versus-geometric multiplicity test, a refusal for neighborhoods that finite precision cannot resolve, enforcement of the representation provenance in which a degeneracy was found, and a biorthogonal response layer with channel-declared residues and sensitivities. It further supplies the spectroscopy-side question that the scalar boundary leaves open, namely when a measured linewidth delivers the amplitude-damping coefficient that a mechanical \(\chi\) requires at all.

The relationship is therefore one of operationalization rather than rediscovery. The earlier work stated the boundaries; this work makes them testable, and makes the cases outside them refusable.

\noindent\textbf{Reader's guide.} A reader interested in the classification engine should read Sections~2.2 to 2.5 (vocabulary, safeguards, response layer) and the synthetic model-class cross-check in Section~2.6. A reader interested in reading a linewidth as a damping ratio should read Sections~5 and 6. The claim-status table in the Supplementary Material records what is established, what is demonstrated, and what is explicitly not claimed.

\section{Method: generator-first stability classification}

\subsection{The unified object}

Every supported input reduces to the spectrum of one generator:
\begin{itemize}
\item Second-order (vibrations, phonons, adsorbate modes): \(M\ddot{q} + C\dot{q} + Kq = 0\), solved as the quadratic eigenvalue problem \citep{tisseur}, \((\lambda^2 M + \lambda C + K)u = 0\), via companion linearization, with the left and right eigenvectors of the response layer following the non-self-adjoint spectral construction of \citet{keldysh}.
\item First-order (microkinetics, master equations): \(\dot{z} = Gz\).
\end{itemize}
For a scalar second-order mode the poles are \(\lambda = -\gamma/2 \pm i\sqrt{\omega_0^2 - \gamma^2/4}\), and the damping ratio \(\chi = \gamma/2\omega_0\) separates the complex-pole (oscillatory, \(\chi<1\)) regime from the real-pole (monotone, \(\chi>1\)) regime, with a repeated root at \(\chi=1\).

\subsection{Classification vocabulary}

ChemSA reports, per mode, four orthogonal properties: stability (stable, marginal, unstable); temporal form (oscillatory, monotone, repeated-root; the term critically damped is reserved for a second-order mechanical reading, where it is licensed); structure (simple, semisimple-at-tolerance, defective-at-tol, higher-order-ep-candidate, near-coalescent-unresolved, unresolved-representation), reported alongside provenance and a scoped interpretation; and mechanical role (well, saddle, soft-mode, node). The bare label \texttt{defective} is not produced for any multiplicity: floating-point evidence, however clean, is tolerance-qualified by construction, and \texttt{defective-at-tol} says so; a candidate of order \(m>2\) that cannot be certified even at that qualification is labelled \texttt{higher-order-ep-candidate} instead (Section~2.3(b)). A mode is not forced into a single label; a system with a saddle and a stable oscillation reports both.

These properties belong to one of three layers, which answer different questions and should not be read as one universal scalar:

\begingroup\small
\noindent\begin{tabular}{@{}p{0.10\linewidth}p{0.42\linewidth}p{0.40\linewidth}@{}}
\toprule
Layer & Object & Question \\
\midrule
A: state & poles, mechanical $\chi$ where licensed, spectral pole-angle $\rho$, temporal and multiplicity class & Where is the mode? \\
B: motion & finite pole shift $(-\Delta\mathrm{Re}\,\lambda, \Delta\mathrm{Im}\,\lambda)$ & What changed between two supplied generators? \\
C: response & left and right eigendata, residues, sensitivity, conditioning & How is it excited, observed, and moved? \\
\bottomrule
\end{tabular}
\endgroup

Layer B is descriptive and path-dependent: it denotes the finite difference between two supplied generator spectra and is retained here as a conceptual comparison, not as a supported software module or a validated general real-system diagnostic. Layer C is predictive and channel-dependent: it maps declared perturbations and drive/readout channels to pole motion and observable response (Section~2.5). Neither layer is another universal scalar alongside $\chi$.

\subsection{The safeguards}

Six safeguards distinguish ChemSA from a naive linewidth-to-\(\chi\) conversion, in two categories: three core scientific safeguards, and three implementation validity gates that make the first three trustworthy in practice.

\textbf{Core scientific safeguard 1: exceptional-point versus semisimple discrimination.} This is the core safeguard, and the one a frequency-and-linewidth reading cannot reproduce. Consider two \(2\times2\) generators with the same repeated eigenvalue \(\lambda = -1\):
\begin{itemize}
\item \(G_1 = \mathrm{diag}(-1,-1) = -I\), a repeated root with a full eigenbasis (every vector is an eigenvector). This is a semisimple degeneracy: the modes are independent and separate smoothly under perturbation. In parameterized Hermitian or normal systems this nondefective case is often called a diabolic point; the more general term semisimple is used here, since a repeated scalar matrix is not necessarily a conical diabolic point in the strict sense.
\item \(G_2 = \left[\begin{smallmatrix}-1 & 1\\ 0 & -1\end{smallmatrix}\right]\), the same repeated root \(-1\) but a single eigenvector (geometric multiplicity 1 less than algebraic multiplicity 2). This is a defective matrix: a genuine exceptional point with square-root perturbation sensitivity.
\end{itemize}
Their eigenvalues are identical; only the eigenvectors differ. ChemSA computes algebraic and geometric multiplicity separately and returns \texttt{semisimple-at-tolerance} for \(G_1\) and \texttt{defective-at-tol} for \(G_2\), demonstrating that it measures eigenvector deficiency, not eigenvalue closeness. Its contribution over a textbook Jordan-form check is extending this to the second-order quadratic eigenvalue problem via companion linearization, with the \(\sqrt{\varepsilon}\) tolerance appropriate to EP2 (algebraic-multiplicity-two) defective-eigenvalue splitting -- an order-\(m\) Jordan block generically splits as \(\varepsilon^{1/m}\) instead, the subject of Section~2.3(b) -- and with physical input contracts that reject non-mechanical generators. The discrimination is verified by randomized regression sweeps over constructed matrices, and separately by one physically parameterized model-class cross-check (Section~2.6); neither alone would suffice. 500 randomized-parameter twin oscillators produced zero false exceptional points, and 500 randomized-parameter exact-critical systems produced zero misses (Supplementary~S3, reproducible from the test suite with a fixed seed). The suite is adversarial: the exact cases that exposed defects during development are retained as permanent regression tests. Because a finite-precision eigensolver cannot supply a structural or symbolic certificate of exact defectiveness, even the exact-entry Jordan example is reported as \texttt{defective-at-tol}; the unqualified \texttt{defective} label is not produced by the numerical classifier.

\textbf{Core scientific safeguard 2: provenance gating.} Defectiveness established from a supplied equation is defectiveness of that representation. A reduced or effective generator may have a genuine exceptional point of its own effective dynamics; what it cannot do, without further analysis, is establish defectiveness of the parent generator it came from, since a repeated root produced by eliminating variables can be an artifact of the elimination. This is an interpretive constraint, not a measurement one: the multiplicities do not change because a flag changed. ChemSA therefore reports three separate fields, structure (the numerical finding, provenance-independent), provenance (full, reduced, or unknown, as declared by the caller), and interpretation (what the structure licenses, provenance-dependent), rather than one overloaded label. An earlier version folded these into a single degeneracy label, so the same defective cluster read differently under full and reduced provenance, which erased the measurement in order to state the caveat. The gate remains a mandatory disclosure and interpretation guard, not an automated truth detector: it cannot catch a mislabelled input, only prevent a silent upgrade from reduced to parent.

\textbf{Core scientific safeguard 3: the constructive scope fence.} The third safeguard is a limit on what the engine may be asked for at all. Local generator architecture classifies a well; it does not determine a rate. Section~3 makes this constructive rather than declarative: two systems with IDENTICAL reactant wells, and therefore identical classification by the engine, differ in activated rate by a factor of $e^{5} \approx 148$ through the global barrier alone. The fence is executable, is run as a validator stage, and classifies both wells through the same public engine the rest of this work uses, so a drift in that engine would break it.

\textbf{Implementation validity gate 1: source gating of numerical inputs.} The provenance gate above governs the representation an equation came from; a second, data-side gate governs whether a number is entitled to appear at all. Every quantity entering a reported calculation is carried as a value paired with its source, and a value with no source is held to be absent: it contributes no number, and any derived quantity that depends on it is returned blocked rather than estimated, inheriting the weakest provenance of its inputs. It is a disclosure and blocking mechanism, not a check that a cited source actually supports its value; that remains a human responsibility. The worked example of Part~III exercises it directly: its one quantitative result, the endpoint width growth, is computed through the gate from the two source-anchored widths, so an illustrative digitized point is structurally unable to contribute to it.

\textbf{Implementation validity gate 2: scalar-\(\chi\) applicability.} A single scalar \(\chi\) is meaningful only when the damping is proportional (\(C\) and \(K\) simultaneously diagonalizable in the mass metric, the Caughey--O'Kelly condition \citep{caughey}). This condition is TESTED by the engine, not accepted on declaration: \texttt{scalar\_reduction\_valid} computes the relative mass-metric commutator of \(M^{-1}C\) and \(M^{-1}K\) and withholds the mechanical \(\chi\), with a stated reason, whenever it is nonzero. When it is not, ChemSA reports not-applicable rather than forcing a number.

\textbf{Implementation validity gate 3: finite-precision equality and certified-domain caution.} A tolerance cluster is not a proof of exact equality, in either direction. Just as a numerically coalescent pair is labelled defective-at-tolerance rather than a certified exceptional point, a repeated eigenvalue with a full eigenbasis is labelled semisimple-at-tolerance, never as exact semisimplicity, which finite precision cannot establish. The reserved label semisimple-exact is used only for multiplicity established analytically or structurally, and the engine never asserts it from a float computation.

\subsection{Scale and representation invariance}

Quadratic eigensolution is performed in mass-normalized coordinates, which removes ordinary coordinate-unit scaling and improves congruence robustness. The certified numerical range is empirical and limited by conditioning: the release reports the mass condition number, applies the tested acceptance gates, and refuses cases outside the calibrated domain. No finite-precision implementation is claimed to be invariant under arbitrarily ill-conditioned congruences; a random non-diagonal congruence was measured to defeat exact-coincidence certification at a mass condition number as low as \(10^2\), well below machine-precision limits, so certification is claimed only within the conditioning range the adversarial suite (well-conditioned, largely diagonal or near-diagonal mass matrices) actually exercises. The governing principle is stated plainly: conditioning can establish that an equality is unresolved; it can never establish that an equality holds. A cluster that shows no deficiency and whose members are not resolvably separated is reported as near-coalescent-unresolved rather than as semisimple.

The certified multiplicity domain is stated and enforced rather than assumed, and it is narrower than two earlier releases of this engine claimed. ChemSA certifies only simple roots and isolated algebraic-multiplicity-two clusters, decided from finite-precision evidence alone: order 2 is provably decidable this way, since two distinct eigenvalues never produce a nonzero rank-one nullity at their mean. NO other verdict is ever certified from floating-point evidence alone. A prior release carved out an exception for members equal to within the linear roundoff radius -- an apparently exact higher-multiplicity coincidence, such as a critically damped twin whose four poles all equal \(-1\) to machine precision -- reasoning that no alternative diagonalizable system could produce that pattern at every achievable precision. That reasoning was WRONG: a triangular matrix with three eigenvalues differing by \(10^{-15}\) and \(2\times10^{-15}\) is genuinely, exactly distinct and diagonalizable, since \(K=1-10^{-15}\) is an ordinary, exactly representable float64 stiffness rather than roundoff noise, yet it fell within that same radius and was certified defective with full confidence. Floating-point proximity, no matter how tight, is never proof of exact mathematical coincidence, so that carve-out has been removed. For a candidate of size \(m>2\), two independent tests are run: a Jordan-chain continuation test (confirming a generalized eigenspace of dimension exactly \(m\) with genuine deficiency) and a Puiseux-fan geometry test (confirming its members are spread the way \(m\) genuine roots of a common defect are). Passing both is NECESSARY evidence for a genuine order-\(m\) defect; it is NOT SUFFICIENT, and is reported only as \texttt{higher-order-ep-candidate}, never certified, REGARDLESS of how small the members' spread is -- a critically damped twin now reports as a candidate like everything else of its size. The reason is structural, not a numerical shortfall: a real, non-roundoff perturbation of a Jordan block, \(J_m(-1)\) with its \((m,1)\) entry offset by an actual \(\varepsilon\), has characteristic polynomial \((\lambda+1)^m=\varepsilon\) and therefore \(m\) genuinely distinct, fully resolvable, diagonalizable roots for every \(\varepsilon\neq0\) -- separated by as much as \(10^{-2}\) at \(m=6\) -- and both tests pass this system identically to a genuine defect, because the Puiseux fan they detect is produced by PERTURBING a Jordan block, and observing that fan does not prove the perturbed matrix remains one. No finite-precision test on a single matrix instance can settle that in general; only a structural or symbolic argument, or parameter continuation establishing that geometric multiplicity genuinely collapses along a swept control parameter, can. A candidate component is refused outright, and returned near-coalescent-unresolved without any label, whenever no plausible local partition can be established at all: a repeated pair not isolated from the nearest external pole by more than its clustering radius, or a component of more than two members that fails either validation test. This boundary is exercised by mandatory crowded-spectrum and order-aware regression suites (\texttt{test\_crowded\_spectrum.py}, \texttt{test\_order\_aware\_ep.py}) wired into the master validator, the latter verified against zero false certifications across perturbation scales from \(10^{-2}\) to \(10^{-12}\), including the triangular-matrix and distinct-scalar-QEP counterexamples that motivated this correction.

\subsection{Biorthogonal response layer}

Because eigenvalue-based readings cannot determine observables (two systems can share identical poles, \(\chi\), and pole motion yet produce different measured spectra), a response layer built from left and right eigenvectors supplies channel-specific residues, structured perturbation sensitivities, and declared-weight condition numbers. For isolated or uniquely partitioned nondefective clusters passing the stated separation and conditioning gates, the response layer reconstructs the supplied-coordinate resolvent, and refuses at defective poles, where the simple-pole residue is invalid rather than merely inaccurate. Crowded neighborhoods whose repeated subset cannot be uniquely isolated are returned near-coalescent-unresolved and excluded from response reconstruction rather than assigned a multiplicity: a candidate component with more than one plausible local partition (for example an exceptional point beside a nearby simple pole) is refused by an isolation and partition-uniqueness gate before any structure or response is reported. This boundary is enforced in code and exercised by mandatory crowded-spectrum regression tests (\texttt{test\_crowded\_spectrum.py}, wired into the master validator). The layer enforces the reciprocal mechanical input contract (\(M\) Hermitian positive definite, \(C\) and \(K\) Hermitian), so the two layers cannot accept different physical domains. Passivity is not assumed: negative damping and negative stiffness remain classifiable as active or unstable. The response layer accepts Hermitian complex coefficients; the classifier path is exercised and tested on real symmetric coefficients, and complex Hermitian input to the classifier is outside the tested domain. A second, physically observable signature of defectiveness follows from this layer: under a generic perturbation of size \(\varepsilon\), eigenvalues at a semisimple degeneracy split linearly in \(\varepsilon\), while at an EP2 (the only order this signature is checked against) they split as \(\sqrt{\varepsilon}\). Computed log-log slopes are 1.00 (semisimple) and 0.50 (defective), corroborating the multiplicity test of Section~2.3(a) through a completely different mechanism.

\subsection{External cross-check: an experimentally realized exceptional-point model class}

The demonstrations in Section~2.3 use constructed matrices, which is the right way to show what eigenvector deficiency means but does not test the engine against anything the world produced. That gap is closed with a system in which an exceptional point has been realized and measured in a laboratory: two linearized spring-coupled pendulums with two real controllable parameters, the length of one pendulum (moving the mass and stiffness matrices together) and a viscous-like damping from an eddy-current brake (moving the damping matrix alone) \citep{even2024}. That work reports, independently, coalescence of both eigenvalues and eigenvectors at the exceptional point, detection of eigenvector coalescence through the condition number, and Puiseux (fractional-exponent) eigenvalue behaviour in its neighbourhood. It also discriminates the point by comparing a pure-exponential fit of its measured response against a Jordan-chain fit, which is the same secular criterion applied in the time domain. Every constant used here is the authors' own calibrated value, identified by them from their experiment and taken at the full precision released in their software \citep{evencode}. The exceptional point is then located by minimizing the eigenvalue gap over the same two real controllable parameters. That search is a local refinement initialized in the neighbourhood of the published coordinates, not an answer-blind global search, and it is reported as such; what makes it informative is that the published six-figure coordinates do not themselves return a degenerate pair. It locates the point at pendulum length \(l_2 = 0.674843\)~m and brake damping \(c_2 = 0.132988\)~N\,m\,s (residual eigenvalue gap \(4.4\times10^{-9}\)), against their published \(0.674831\)~m and \(0.133028\)~N\,m\,s: agreement to \(1.2\times10^{-5}\)~m and \(4.0\times10^{-5}\)~N\,m\,s. The eigenvalue there, \(-0.023143 - 3.271192i\), matches their reported \(-0.023155 \pm 3.271194i\) to \(2.5\times10^{-6}\)~rad\,s\(^{-1}\) in frequency. Evaluating instead at their published six-figure coordinates returns two simple poles separated by about \(10^{-3}\), which is not a disagreement but the codimension-two Puiseux sensitivity the source itself emphasizes: six figures do not place a system on an exceptional point, so it must be relocated in the same parameters. The reported signatures reproduce: the structure verdict is defective-at-tolerance; algebraic multiplicity 2 exceeds geometric multiplicity 1; the declared-weight condition number diverges on approach and is refused at the exceptional point itself. The pendulum validates the mathematics (a quadratic eigenvalue problem solved by companion linearization); a phonon application additionally requires the reciprocal mechanical input contract of Section~2.5, which Hermitian force-constant and damping matrices satisfy. The authors' measured time traces are not ingested, so this reproduces their calibrated model rather than their raw data, and it is not an independent experimental validation. What it establishes is narrower and sufficient: that the machinery fires on a physically parameterized, laboratory-realizable model class located by search, rather than only on a constructed Jordan block.

\section{The scope boundary (constructive)}

Two systems are constructed with identical local reactant wells (same \(M\), \(C\), \(K\)) but different global barriers. ChemSA returns the same class and identical poles for both. Because the two systems share the well, friction, and barrier frequency, the Kramers--Grote--Hynes prefactor cancels between them exactly, leaving \(k_A/k_B = e^{(\Delta G_B - \Delta G_A)/k_BT} = e^{5} \approx 148\): no absolute Kramers rate is evaluated, only this prefactor-free ratio, since the friction used is deeply underdamped and outside the moderate-to-high-friction regime the prefactor itself is derived for. Identical local stability architecture, rates differing by 148-fold: local classification cannot determine rate. ChemSA therefore classifies mode stability and does not infer reaction rate from local architecture alone. A transition state (a stiffness matrix with a negative eigenvalue) is always classified unstable, never as a stable-well mode, for all frictions in the scalar inverted-harmonic model tested.

\subsection*{Where the coordinate does not exist}

The qualifying rule of Section~1.3 restricts \(\chi\) to a stable second-order response. There is a sharper statement available on the other side of the surface, and it is algebraic rather than evidential. For a stable mode the characteristic discriminant is \(\gamma^2/4 - \omega_0^2\), which vanishes at \(\gamma = 2\omega_0\); that vanishing is the boundary \(\chi = 1\). For an inverted mode, a transition state with negative curvature along the reaction coordinate,
\[
\ddot q + \gamma \dot q - \omega_b^2 q = 0, \qquad \text{discriminant} = \frac{\gamma^2}{4} + \omega_b^2 ,
\]
and this cannot vanish for real \(\gamma\) and nonzero \(\omega_b\). The roots remain real and distinct at every friction, one of them always positive. A dimensionless friction remains perfectly well defined at a barrier. Writing it \(\alpha_b = \gamma/2\omega_b\), the normalized barrier friction, it parameterizes \(\lambda_+\) and \(\kappa_b\) in the scalar passive model and governs how much of an incoming flux is transmitted. It is deliberately not called a mechanical damping ratio, because that name carries a critical value and this quantity has none. What is absent is the boundary: the roots never coalesce, so nothing changes character at \(\alpha_b = 1\) and the coordinate carries no landmark on the side of the surface where reactions occur.

This is the reason a damping ratio cannot carry a rate, stated without a counterexample. The quantity that does govern transmission through a barrier is the positive reactive pole \(\lambda_+ = -\gamma/2 + \sqrt{\gamma^2/4 + \omega_b^2}\), or equivalently the transmission factor \(\kappa_b = \lambda_+/\omega_b\) of Kramers and of Grote and Hynes \citep{hanggi1990}, which is monotonically decreasing and so has no maximum either. Locating exceptional points in parametric quadratic eigenproblems is itself established: continuation methods and the open-source \texttt{EasterEig} implementation compute defective eigenvalues and their Puiseux expansions directly \citep{eastereig}. The contribution here is not the location of such points but the conservative workflow around them, namely classification of a supplied representation, refusal of neighbourhoods that finite precision cannot resolve, enforcement of representation provenance, and channel-declared response geometry. The engine already enforces this: a stiffness matrix with a negative eigenvalue is classified as a saddle and \(\chi\) is withheld for the reactive coordinate at every friction tested, with the reactive pole reported in its place. Stable transverse modes at the same saddle retain their own damping ratios, which is why the statement is about the reactive direction rather than about the whole generator. Even that pole is only a prefactor, since the activation exponent is not a property of the local generator, which is the same boundary the constructive counterexample of Section~3 establishes by a different route.

\section{The damping-ratio coordinate and its qualifying rule}

\subsection{The qualifying rule}

A cell qualifies for the \(\chi\) axis only if its reported damping is the amplitude-damping \(\gamma\) of a damped-harmonic-oscillator (pole) response, the same quantity as \(\chi = \gamma/2\omega_0\). Transition linewidths, total-coherence (\(T_2\)) and pure-dephasing (\(T_\phi\), equivalently \(T_2^*\)) rates, and pressure-broadening coefficients are different physical quantities and are excluded however cleanly measured. This is why quantities that are not amplitude damping cannot share the coordinate: a dephasing rate and an oscillator damping are not silently placed on the same coordinate. The single worked example of Part~III is traced to its primary source, with the damping convention confirmed against the paper.

\subsection{The one-axis observation}

The regimes are one axis, not separate phenomena: under (\(\chi<1\)), critical (\(\chi=1\)), over (\(\chi>1\)). The critical point is a boundary systems cross, not a state a real system rests at, since \(\chi=1\) is where oscillation gives way to monotone decay, a threshold crossed rather than occupied.

A concrete over-damped anchor: the M-tilt phonon mode of CsPbBr\(_3\) at 350~K, from an atomistic-simulation study of the cubic-to-tetragonal transition \citep{fransson}, whose own equation for mode-coordinate autocorrelation gives two relaxation times at that temperature, \(\tau_S=0.31\)~ps and \(\tau_L=5.22\)~ps. Inverting the source's own relation (\(\tau_{S,L} = \tau/(1\mp\sqrt{1-(\omega_0\tau)^2})\), \(\tau=2/\Gamma\)) gives \(\Gamma = 1/\tau_S+1/\tau_L = 3.417\)~ps\(^{-1}\) and \(\omega_0 = 1/\sqrt{\tau_S\tau_L} = 0.786\)~ps\(^{-1}\), so \(\chi=\Gamma/2\omega_0=2.174\): genuinely over-damped, and the classifier returns monotone, stable, \(\chi>1\) for it with no oscillatory component. The round-trip through the source's own equation reproduces both published times to \(5\times10^{-15}\), and the regime survives a hypothetical 10\% perturbation on either time (Supplementary, \texttt{verify\_cspbbr3\_cell.py}). This is a simulated, not measured, linewidth, and is reported as such: the pass table above admits only directly measured population lifetimes, so this entry stays here rather than in it.

\subsection{The hinge}

The qualifying rule of Section~4.1 excludes dephasing, pressure broadening, and inhomogeneous width from the $\chi$ axis on the ground that they report different physical quantities. That exclusion is correct, and it raises a sharper question than it answers: given a measured linewidth, under what conditions may it legitimately be read as a mechanical $\chi$ at all? Part~III answers that question. The answer is not a second and third coordinate on which those excluded quantities are plotted (that would repeat the error the qualifying rule exists to prevent, placing physically distinct objects on a common-looking axis); it is a hierarchy of gates through which a linewidth must pass before Section~4.1's coordinate applies to it. The engine classifies whether a degeneracy is defective; the hierarchy classifies whether a linewidth may become a $\chi$. The two are the same epistemic move applied to two different objects, and neither reads its answer from the surface number.

\section{A three-step procedure for reading a linewidth as \texorpdfstring{$\chi$}{chi}}

\subsection{What a linewidth is made of}

A measured spectral linewidth is, in general, not one quantity. For a mode of frequency $\omega_0$ the homogeneous width decomposes as
\begin{equation}
\widetilde{\Gamma}_{\mathrm{hom}} = \frac{1}{\pi c T_2} = \frac{1}{2\pi c T_1} + \frac{1}{\pi c T_\phi} + \widetilde{\Gamma}_{\mathrm{or}},
\end{equation}
where $T_1$ is the population (energy) relaxation time, $T_\phi$ is the pure-dephasing time and $\widetilde{\Gamma}_{\mathrm{or}}$ is the orientational contribution. A tilde marks a width in WAVENUMBERS (cm$^{-1}$), the unit in which every width in this work is quoted and compared; the same decomposition written without $c$ is a set of rates in s$^{-1}$, and the factor $c$ between them is exactly the kind of conversion slip the survey cited above reports as prevalent, so it is carried explicitly rather than left implicit. Only the first term is a mechanical damping coefficient, and the absorption line may carry, in addition, a static inhomogeneous spread of centre frequencies across sites. The mechanical damping coefficient is $\gamma = 1/T_1$, and the coherence-amplitude contribution to the homogeneous width is $\gamma/2 = 1/2T_1$; only this lifetime part is an amplitude decay, and only it becomes a mechanical damping ratio through $\chi = \gamma/2\omega_0$. Reading a total width as $\chi$ without separating these contributions places a dephasing rate or a static spread onto the mechanical coordinate, which is precisely the failure the qualifying rule forbids. Three conditions must therefore be checked, in order, before a linewidth yields a mechanical $\chi$. Each is a single physical ratio, and each is a gate rather than a coordinate.

\subsection{Step 1: resolvability}

A per-mode analysis presupposes that modes are resolved. Two ratios govern this. For a single line, the static inhomogeneous spread relative to the homogeneous width,
\begin{equation}
R_{\mathrm{abs/hom}} = \frac{\Gamma_{\mathrm{abs}}}{\Gamma_{\mathrm{hom}}},
\end{equation}
states whether the observed width reports homogeneous dynamics ($R_{\mathrm{abs/hom}} \lesssim 1$) or is dominated by a static distribution that an echo must remove before $T_1$ and $T_\phi$ can be read. For a pair of nearby modes separated by $\Delta$, the pair-resolvability ratio
\begin{equation}
R_{\mathrm{pair}} = \frac{\Delta}{\Gamma_{\mathrm{total}}}
\end{equation}
states whether the two lines are separable at all. Both ratios depend on the assumed line shape and on a stated resolution criterion; they are not universal constants. In the illustrative case of two equal Gaussian lines a merger occurs near $\Delta \approx 2\sigma$, but this is a line-shape-specific and criterion-specific boundary, not a universal critical value, and it is reported as such. Step 1 passes when a mode is resolvable and its homogeneous width is separable from static heterogeneity.

\subsection{Step 2: interpretability}

Once a homogeneous width is in hand, it must be separated into its lifetime and pure-dephasing parts before it can be read as amplitude damping. The governing ratio is
\begin{equation}
R_\phi = \frac{1/T_\phi}{1/2T_1},
\end{equation}
the pure-dephasing rate relative to the amplitude-decay rate. Its role is to quantify a bias, not to gate recoverability. If the total homogeneous width is read as though it were lifetime damping, the inferred damping ratio is too large by exactly the factor
\begin{equation}
\frac{\chi_{\mathrm{naive}}}{\chi_{\mathrm{true}}} = 1 + R_\phi + R_{\mathrm{or}},
\end{equation}
where $R_{\mathrm{or}} = \widetilde{\Gamma}_{\mathrm{or}}/(1/2\pi c T_1)$ is the orientational contribution relative to the lifetime contribution. The homogeneous width carries three terms, not two \citep{tokmakoff1995},
\[
\widetilde{\Gamma}_{\mathrm{hom}} = \frac{1}{\pi c T_2} = \frac{1}{\pi c T_\phi} + \frac{1}{2\pi c T_1} + \widetilde{\Gamma}_{\mathrm{or}},
\]
and only the lifetime term is a mechanical damping coefficient. The two-term reduction used below is licensed only when $\widetilde{\Gamma}_{\mathrm{or}}$ is independently measured, independently bounded, or negligible. It is not automatically negligible: for a degenerate transition the orientational contribution is not physical rotation of the molecule but time evolution of the coefficients of the initially excited superposition, which can remain fast even where the solvent viscosity changes by many orders of magnitude, and it has been reported to contribute significantly to the homogeneous width at intermediate temperatures while never dominating it \citep{tokmakoff1995}. The worked example of Section~6 uses a nondegenerate transition, which removes the superposition-evolution channel. Ordinary rotational diffusion of the transition dipole is not thereby removed, and it is treated separately. For the worked example the source states that orientational relaxation does not occur on the echo timescale because of the high solvent viscosity, and reports that this was checked independently by polarization-selective pump-probe measurement \citep{rectorfayer}. The two-term reduction used there is therefore supported by evidence on both counts, by symmetry for the superposition channel and by measurement for the rotational one, rather than assumed. That assumption is stated here rather than absorbed silently, and any application in which the orientational correlation time approaches the vibrational lifetime requires \(\Gamma_{\mathrm{or}}\) to be measured or bounded before the reduction is used. Setting $R_{\mathrm{or}} = 0$,
so $R_\phi$ is the relative bias of the naive width reading: $R_\phi = 0.1$ is a ten percent overestimate, $R_\phi = 1$ a factor of two, $R_\phi = 10$ an elevenfold overestimate. The value $R_\phi = 1$ is the equal-contribution crossover, where lifetime and pure dephasing contribute equally to the width; it is not a critical point and not an identifiability transition, and nothing bifurcates there. Recoverability of $\chi$ is a separate matter that $R_\phi$ does not govern: when $T_1$ is measured independently (as by an infrared pump-probe alongside a vibrational echo for $T_2$), $\chi$ is recoverable at every $R_\phi$, including $R_\phi \gg 1$; when only the total width is available, $\chi$ is not uniquely determined for any unknown nonzero dephasing, including $R_\phi < 1$, unless the dephasing contribution is independently bounded. A system for which only the total width is known therefore yields only the weak upper bound $0 \le \chi_{\mathrm{true}} \le \chi_{\mathrm{naive}}$ (from the nonnegativity of the two contributions) unless the dephasing is independently bounded; no point value of $\chi$ is warranted in that case.

\subsection{Step 3: mechanical placement}

The mechanical coordinate is applied only after a mode-specific homogeneous response has been isolated (Step~1) and its lifetime contribution independently measured or bounded against BOTH pure dephasing and orientational relaxation (Step~2); $R_\phi$ and $R_{\mathrm{or}}$ are diagnostics within Step~2, not pass/fail thresholds. It is then computed from the lifetime and the frequency alone,
\begin{equation}
\chi = \frac{1/T_1}{2\omega_0},
\end{equation}
which is the coordinate of Part~II with its input now certified to be an amplitude-decay rate rather than a dephasing rate or a static spread. This ordering is the content of Part~III: the single coordinate of Part~II is not generalized into three coordinates, it is protected by two gates that state when it is meaningful. A degeneracy is not self-describing, and Part~I resolves that by testing eigenvector deficiency rather than eigenvalue closeness; a linewidth is not self-describing, and Part~III resolves that by testing resolvability and lifetime-versus-dephasing content rather than reading the surface width.

\section{One worked trajectory: Rh(CO)\texorpdfstring{$_2$}{2}acac}

The hierarchy is exercised on a system in which $T_1$ and $T_2$ were measured independently on the same transition, so that the lifetime and dephasing contributions are separately grounded rather than assumed. For the asymmetric CO stretch of Rh(CO)$_2$acac (2010~cm$^{-1}$, A$_1$ of C$_{2v}$, nondegenerate) in dibutyl phthalate over 3.4--250~K \citep{rectorfayer}, a vibrational echo measures $T_2$ and an infrared pump-probe measures $T_1$ (reported as $2T_1$), decomposed through the source's Eq.~1, $\Gamma = 1/\pi T_2 = 1/\pi T_2^* + 1/2\pi T_1$, which is written there as a rate; dividing by $c$ gives the wavenumber widths quoted below, $\widetilde{\Gamma} = 1/\pi c T_2$. Step~1 is non-trivial here: at 3.4~K the homogeneous width is 0.11~cm$^{-1}$ inside an absorption envelope of $\approx$15~cm$^{-1}$, a ratio of order $10^2$, so the raw absorption is a static-inhomogeneous envelope and cannot be read as a damping coefficient until the echo isolates the homogeneous line. The source states that $2T_1$ has a very mild (linear) temperature dependence while $T_2$ falls steeply, and that the lifetime at each temperature was taken from the linear fit to all the $T_1$ data.

Figure~\ref{fig:hierarchy} carries the transition through the three steps. Two points are quantitative, both source-text-anchored: at 3.4~K the homogeneous width is 0.11~cm$^{-1}$ ($T_2 = 95.2$~ps) and at 250~K the source-anchored width is 1.5~cm$^{-1}$, from which $T_2 = 7.08$~ps is author-derived rather than stated. At both, $\chi = (1/T_1)/2\omega_0$ from the pump-probe lifetime is near $3\times10^{-5}$, deeply underdamped. The quantitative result is the endpoint contrast: the homogeneous width grows by a factor of about fourteen as pure dephasing comes to dominate, while the mechanical $\chi$ barely moves, so reading the 250~K total width as damping overestimates $\chi$ by roughly twelvefold. The intermediate echo and pump-probe points in Figure~\ref{fig:hierarchy} are by-eye readings of the source figure and are shown as an illustrative trajectory only; no crossover temperature or intermediate bias factor is reported as a result, because a by-eye digitization does not support one. This is one worked real example, not a catalogue, and it makes the central point precisely: coherence can be lost rapidly while the mechanical mode remains deeply underdamped, so the total linewidth and the mechanical damping ratio are different quantities and only an independent lifetime measurement yields the latter.

The two halves of this work meet at this point rather than remaining parallel. At each temperature the fitted lifetime defines a scalar second-order generator, $M = 1$, $C = \gamma = 1/T_1$, $K = \omega_0^2$, which is fed through the classifier of Part~I (\texttt{classify\_modes}). The engine returns, independently of the linewidth calculation, a stable and oscillatory simple mode at every temperature, and its structural damping coordinate agrees with the directly computed $\chi = \gamma/2\omega_0$ to machine precision across the full range. The linewidth procedure and the degeneracy engine are therefore one pipeline: a measured coherence trajectory is reduced to a certified lifetime, the lifetime builds a generator, and the same engine that distinguishes semisimple from defective multiplicity classifies that generator and confirms its placement. The agreement is verified as a regression test in the shipped suite.

\begin{figure}[t]
\centering
\includegraphics[width=\textwidth]{linewidth_hierarchy.png}
\caption{The three-step procedure carried through one real transition, the asymmetric CO stretch of Rh(CO)$_2$acac (Rector \& Fayer 1998). Filled squares in the centre panel mark the two source-text-anchored homogeneous-width endpoints and are the only quantitative plotted values. Open symbols and dashed curves, including every lifetime-derived \(\chi\) and every bias value in the right panel, are source-informed illustrations built from by-eye transcriptions of the published pump-probe figure. Left (Step~1, homogeneous isolation at 3.4~K): the observed absorption envelope ($\approx$15~cm$^{-1}$) is about $1.3\times10^2$ times broader than the echo-isolated homogeneous width ($\approx$0.11~cm$^{-1}$), so the observed absorption width cannot be used as a mechanical damping coefficient. Centre (Step~2, decomposition): at the two endpoints the homogeneous width grows from 0.11 to 1.5~cm$^{-1}$ while the lifetime part stays nearly constant, so pure dephasing accounts for the growth. Right (Step~3, mechanical placement): $\chi$ from the lifetime stays near $3\times10^{-5}$ (deeply underdamped) at both endpoints; the dotted curve is the biased estimate from reading the total homogeneous width as damping, about twelvefold too large at 250~K.}
\label{fig:hierarchy}
\end{figure}

\subsection*{The interior of the axis}

The entries that follow are chosen because they stress the criterion. It is worth stating first what the ordinary case looks like, since a procedure demonstrated only on its boundary invites the impression that it mostly refuses. Thirty systems, spanning the gas phase, ordinary liquids, water, and eleven solid-state phonon modes, all satisfy the qualifying rule and all return a mechanical damping ratio with no refusal. Every value is a directly reported amplitude-damping quantity: optoacoustic vibration-translation relaxation in the gas phase \citep{cannemeijer1973,cannemeijer1974,ch4vt,cottrell1966}, infrared pump-probe population lifetimes in the condensed phase \citep{tokmakoff1994,steinel2004,owrutsky1991}, and a classical damped-harmonic-oscillator (Lorentz-oscillator) dielectric fit for eleven Raman-inactive or infrared-active phonon modes of two solids: three modes of the perovskite KTaO$_3$ \citep{mori2026} and eight room-temperature modes (five ordinary-ray, three extraordinary-ray) of the mineral calcite, CaCO$_3$ \citep{posch2007}. Both solid-phase sources report the amplitude-damping constant directly in a numeric table, not as a lineshape width and not read from a figure. The Rh(CO)$_2$acac system of Section~6 is deliberately NOT among them: its lifetime is not source-tabulated, and an entry whose lifetime would have to be inferred from a width is exactly what this table must not contain. No spectral width is converted.

\begin{table}[htbp]
\centering
\noindent\footnotesize\setlength{\tabcolsep}{3pt}\begin{tabular}{@{}p{0.17\linewidth}p{0.05\linewidth}p{0.08\linewidth}p{0.10\linewidth}p{0.185\linewidth}p{0.075\linewidth}p{0.16\linewidth}@{}}
\toprule
System & Phase & \(\nu_0\)/cm\(^{-1}\) & \(\tau\)\(^{\dagger}\) & \(\chi\) & Tier & Source \\
\midrule
CO$_2$ \(\nu_2\) bend, pure CO$_2$ & gas & 667.0 & 6.70 & \(5.94\!\times\!10^{-10}\) & abs. & \citep{cannemeijer1973} \\
CO$_2$ \(\nu_2\) bend, CO$_2$-Ne & gas & 667.0 & 10.20 & \(3.90\!\times\!10^{-10}\) & abs. & \citep{cannemeijer1973} \\
CO$_2$ \(\nu_2\) bend, CO$_2$-Ar & gas & 667.0 & 54.00 & \(7.37\!\times\!10^{-11}\) & abs. & \citep{cannemeijer1973} \\
CO$_2$ \(\nu_2\) bend, CO$_2$-Kr & gas & 667.0 & 200.00 & \(1.99\!\times\!10^{-11}\) & abs. & \citep{cannemeijer1973} \\
CO$_2$ \(\nu_2\) bend, CO$_2$-N$_2$ & gas & 667.0 & 12.80 & \(3.11\!\times\!10^{-10}\) & abs. & \citep{cannemeijer1974} \\
CO$_2$ \(\nu_2\) bend, CO$_2$-O$_2$ & gas & 667.0 & 8.80 & \(4.52\!\times\!10^{-10}\) & abs. & \citep{cannemeijer1974} \\
CH$_4$ \(\nu_4\), pure CH$_4$ & gas & 1306.0 & 1.65 & \(1.23\!\times\!10^{-9}\) & abs. & \citep{ch4vt} \\
CH$_4$ \(\nu_4\), CH$_4$-Ar & gas & 1306.0 & 12.00 & \(1.69\!\times\!10^{-10}\) & abs. & \citep{ch4vt} \\
W(CO)$_6$ in CCl$_4$, 295 K & liquid & 1980.0 & 700.00 & \(1.92\!\times\!10^{-6}\) & text & \citep{tokmakoff1994} \\
W(CO)$_6$ in CHCl$_3$, 295 K & liquid & 1980.0 & 340.00 & \(3.94\!\times\!10^{-6}\) & text & \citep{tokmakoff1994} \\
Cr(CO)$_6$ in CCl$_4$, 295 K & liquid & 1980.0 & 370.00 & \(3.62\!\times\!10^{-6}\) & text & \citep{tokmakoff1994} \\
Cr(CO)$_6$ in CHCl$_3$, 295 K & liquid & 1980.0 & 260.00 & \(5.16\!\times\!10^{-6}\) & text & \citep{tokmakoff1994} \\
OD stretch, HOD in H$_2$O & water & 2500.0 & 1.45 & \(7.32\!\times\!10^{-4}\) & text & \citep{steinel2004} \\
CH$_4$ \(\nu_4\), CH$_4$-He & gas & 1306.0 & 1.97 & \(1.03\!\times\!10^{-9}\) & abs. & \citep{ch4vt} \\
CH$_4$ \(\nu_4\), CH$_4$-Ne & gas & 1306.0 & 8.60 & \(2.36\!\times\!10^{-10}\) & abs. & \citep{ch4vt} \\
CH$_4$ \(\nu_3\), pure CH$_4$ & gas & 3019.0 & 1.60 & \(5.50\!\times\!10^{-10}\) & abs. & \citep{cottrell1966} \\
N$_2$O \(\nu_3\), pure N$_2$O & gas & 2224.0 & 1.60 & \(7.46\!\times\!10^{-10}\) & abs. & \citep{cottrell1966} \\
COS \(\nu_3\), pure COS & gas & 2062.0 & 0.70 & \(1.84\!\times\!10^{-9}\) & abs. & \citep{cottrell1966} \\
N$_3^-$ antisym stretch, protic & liquid & 2048.0 & 3.00 & \(4.32\!\times\!10^{-4}\) & abs. & \citep{owrutsky1991} \\
KTaO$_3$ TO1 soft mode & solid & 79.4 & 244.6 fs & \(1.37\!\times\!10^{-1}\) & text & \citep{mori2026} \\
KTaO$_3$ TO$_2$ mode & solid & 197.6 & 672.0 fs & \(2.00\!\times\!10^{-2}\) & text & \citep{mori2026} \\
KTaO$_3$ TO4 mode & solid & 547.0 & 163.9 fs & \(2.96\!\times\!10^{-2}\) & text & \citep{mori2026} \\
Calcite OR mode 1 & solid & 1407.4 & 358.7 fs & \(5.26\!\times\!10^{-3}\) & text & \citep{posch2007} \\
Calcite OR mode 2 & solid & 712.0 & 1327.2 fs & \(2.81\!\times\!10^{-3}\) & text & \citep{posch2007} \\
Calcite OR mode 3, 300K & solid & 297.0 & 338.1 fs & \(2.64\!\times\!10^{-2}\) & text & \citep{posch2007} \\
Calcite OR mode 4, 300K & solid & 222.8 & 482.6 fs & \(2.47\!\times\!10^{-2}\) & text & \citep{posch2007} \\
Calcite OR mode 5, 300K & solid & 104.7 & 884.8 fs & \(2.87\!\times\!10^{-2}\) & text & \citep{posch2007} \\
Calcite ER mode 1 & solid & 872.0 & 2540.1 fs & \(1.20\!\times\!10^{-3}\) & text & \citep{posch2007} \\
Calcite ER mode 2, 300K & solid & 306.7 & 314.1 fs & \(2.76\!\times\!10^{-2}\) & text & \citep{posch2007} \\
Calcite ER mode 3, 300K & solid & 94.3 & 747.7 fs & \(3.76\!\times\!10^{-2}\) & text & \citep{posch2007} \\
\bottomrule
\end{tabular}
\caption{The thirty systems that pass the qualifying rule directly, with no refusal. \(^{\dagger}\)\(\tau\) is the reported population lifetime for the gas and condensed-phase entries; for the eleven solid-phase entries, the source reports an amplitude-damping constant \(\gamma\) directly from a classical damped-harmonic-oscillator (Lorentz-oscillator) dielectric fit, and \(\tau\) shown here is the equivalent \(1/(2\pi c\gamma)\), carrying no new information. \(\chi = \gamma/2\omega_0\) is computed from it through the same engine used throughout this work. Tier records whether the source states the value in its abstract (\emph{abs.}) or its running text/table (Section~2.3); no entry here is a figure reading. The Rh(CO)\(_2\)acac system of Section~6 is deliberately absent: its lifetime is not source-tabulated.}
\label{tab:passtable}
\end{table}

The resulting coordinate spans nearly ten decades, from \(2\times10^{-11}\) for the CO\(_2\) bending mode relaxed by krypton to \(0.137\) for the lowest-frequency (soft) phonon mode of solid KTaO\(_3\). The gas entries fall below the liquid entries, which fall below the single water entry, which falls below the three solid-phase entries; no claim of coupling strength beyond the values themselves is intended, only that this is the order the computed \(\chi\) happens to fall in for the systems this table carries. The solid-phase entries are markedly closer to the critical boundary than anything else in the table, still comfortably underdamped but by a smaller margin; nothing in the qualifying rule treats that as more or less admissible than the gas-phase entries seven orders of magnitude further from it. Every entry is far below the critical boundary, and for the gas phase the size of that gap can be stated in closed form. Vibration-translation relaxation obeys \(\tau(p) = \tau_{\mathrm{atm}}/p\), so \(\chi\) is pressure tunable and reaching \(\chi = 1\) for the CO\(_2\) bending mode would require about \(1.7\times10^{9}\)~atm, some 470 times the pressure at the centre of the Earth. That figure is a formal linear extrapolation of the dilute-gas law and not a pressure at which such a system would in fact become critically damped: \(\tau \propto 1/p\) holds while binary collisions dominate, and the medium at that pressure is not a dilute gas. What the extrapolation establishes is how far outside its own regime the law must be carried before the boundary is approached, and the conclusion survives six decades of assumed relaxation time. The absent near-boundary gas entry is therefore a statement about where the axis has structure rather than about the completeness of a literature search: \(\chi \sim 1\) requires damping comparable to the restoring frequency, and a dilute gas is defined by molecules mostly not interacting.

\section{The procedure as an admission criterion: several published systems}

Section~6 carried one system through all three steps. The procedure is not a
single reconstruction, however, but a decision rule, and its usefulness is
better shown by what it decides across systems that were measured for other
reasons. Table~\ref{tab:admission} records how far each of several published
metal-carbonyl systems gets, and what stops it. Every input is a value stated in
the running text of its source; none is read off a figure. Values a source gives
only as similar to another number are entered as illustrative, and the source
gate of Section~2.3 then blocks any result derived from them from being quoted
as quantitative, without the authors having to remember to do so.

\begin{table}[htbp]
\centering\small
\noindent\begin{tabular}{@{}p{0.19\linewidth}p{0.09\linewidth}p{0.17\linewidth}p{0.08\linewidth}p{0.09\linewidth}p{0.22\linewidth}@{}}
\toprule
System & \(T\) & Verdict & \(R\) & Over & Limiting evidence \\
\midrule
W(CO)$_6$ / 2-MP & 250--300~K & Step 1 passed & 1 & n/a & no text \(T_1\) at this \(T\) \\
W(CO)$_6$ / DBP & 300~K & Step 1 passed & 26 & n/a & no text \(T_1\) at this \(T\) \\
W(CO)$_6$ / CCl$_4$ & 295~K & Step 1 refused & n/a & 1300 & no echo \\
W(CO)$_6$ / CHCl$_3$ & 295~K & Step 1 refused & n/a & 1200 & no echo \\
Cr(CO)$_6$ / CCl$_4$ & 295~K & Step 1 refused & n/a & 700\(^{*}\) & width inherited as similar \\
Cr(CO)$_6$ / CHCl$_3$ & 295~K & Step 1 refused & n/a & 930\(^{*}\) & width inherited as similar \\
Rh(CO)$_2$acac / DBP & 3.4, 250~K & Steps 1--3 & 136 (3.4~K) & n/a & reference entry (Section~6) \\
\bottomrule
\end{tabular}
\caption{The three-step procedure applied as an admission criterion. Verdict records how far each published system gets and what stops it. \(R\) is the ratio of absorption to echo-isolated homogeneous width; Over is the factor by which the absorption width exceeds the lifetime-limited width, a bound on the total overestimate rather than a decomposition. \(^{*}\)Not reportable: an illustrative input, blocked by the source gate.}
\label{tab:admission}
\end{table}

The two entries at the top are the substantive result, and they concern one
molecule and one mode. The asymmetric CO stretch of W(CO)$_6$ near 1980~cm$^{-1}$
in 2-methylpentane above 250~K has an absorption line that is homogeneously
broadened: the echo signal becomes a free induction decay whose width matches
the measured absorption width \citep{tokmakoff1995}, so \(R_{\mathrm{abs/hom}} = 1\)
and reading that width as a dynamic quantity is correct. The same mode of the
same molecule at 300~K in dibutylphthalate has a homogeneous width near
1~cm$^{-1}$ against an absorption width of 26~cm$^{-1}$ \citep{tokmakoff1995}, so
the absorption width overstates the dynamic width by a factor of about 26 and
must not be used. Solvent alone moves the verdict. Whether a measured linewidth
may be read as dynamics is therefore a property of the system rather than of the
chromophore, and no inspection of the spectrum alone distinguishes the two
cases: the echo measurement is what separates them.

The lower entries record the more common situation, in which no echo exists.
For those, Step~1 cannot be satisfied and the procedure returns a bound rather
than a value. The bound is not weak. Comparing the stated absorption width with
the width implied by the stated lifetime, \(\Gamma_{\mathrm{life}} = 1/2\pi c T_1\),
gives overshoot factors near \(10^3\) for the room-temperature carbonyl solutions
\citep{tokmakoff1994}. Reading such a linewidth as a mechanical damping
coefficient would overestimate \(\chi\) by about three orders of magnitude. That
overshoot lumps inhomogeneous broadening together with pure dephasing and
orientational relaxation, and is therefore a bound on the total overestimate,
not a decomposition of it; separating the three requires the echo the procedure
asks for.

\subsection*{A fourth verdict: when the decomposition itself is not well posed}

The entries above are all metal carbonyls, and all of their refusals have the
same character: a measurement is missing. Applying the procedure outside that
family produces a refusal of a different kind, and the distinction matters
because it is a statement about the system rather than about the data.

For the OH stretch of liquid water the vibrational lifetime is not a single
number. Frequency-resolved relaxation measurements give about 250~fs at the red
edge of the band rising to about 550~fs at the blue edge, with a comparable
factor-of-two variation reported for HOD in D$_2$O \citep{water2015}. Step~2 requires one lifetime in order to define a single
scalar mechanical mode. The local physics is not in question: each hydrogen-bond
environment relaxes perfectly well. What fails is the reduction. A macroscopic
band that represents a dispersed ensemble of local environments is not one
oscillator, so mapping the whole band onto a single scalar \((M, C, K)\) is
ill posed, and there is no single macroscopic mechanical coordinate on which to
place it. 

The OD stretch of HOD in H$_2$O separates the two failures cleanly. Its excited
state decay is wavelength independent with a decay time of 1.45~ps \citep{steinel2004}, so Step~2
would be well posed. Step~1 is not: the source reports that the dynamic line width, its
measure of spectral diffusion, reaches 88\% of the total absorption line width by 0.5~ps,
95\% by 1.0~ps, and 98\% by 1.5~ps, so a hydroxyl stretch samples essentially every
environment within one vibrational lifetime. A static inhomogeneous distribution separable
by an echo does not exist at the relevant delay. The decomposition presupposes a
separation of timescales that this system does not supply.

Neither case is closed by simply supplying the missing number: what fails is the static two-component decomposition of Step~1 on the lifetime timescale, and a richer time-dependent treatment, not a further measurement of the same kind, is what the system requires.
Both are systems for which the scalar reduction the
procedure asks for does not have a well-formed answer, and reporting that is the
correct output. The refusal is a statement about the reduction, not a claim that
the underlying dynamics are ill defined. The distinction
between a missing measurement and an ill-posed decomposition is not visible in a
spectrum, and a procedure that returned a number in either case would be wrong
in a way no consistency check would catch.

One further observation from the same source sharpens the point, and it is
worth stating carefully because the obvious reading of it is wrong. For
W(CO)$_6$ in CHCl$_3$ the vibrational lifetime lengthens by about nine percent
between the melting and boiling points, from 322 to 350~ps, while the
absorption linewidth narrows by about nine percent over the same range, from
20.5 to 18.6~cm$^{-1}$ \citep{tokmakoff1994}. A lifetime and a width moving in
opposite senses looks at first like a contradiction, but it is not: since
$\widetilde{\Gamma}_{\mathrm{life}} = 1/2\pi cT_1$, a lengthening lifetime
predicts a narrowing lifetime width, and the two trends are therefore
directionally CONSISTENT ($-8.0$ percent predicted against $-9.3$ percent
observed). No direction test is available here, and none is claimed. What
separates the two quantities is MAGNITUDE. The lifetime-limited widths implied
by those lifetimes are 0.0165 and 0.0152~cm$^{-1}$, so the measured absorption
line is about $1.2\times10^{3}$ times broader at both ends of the range and is
dominated throughout by contributions that are not lifetime. A near-agreement
in trend across two observables separated by three orders of magnitude in size
is exactly the coincidence that makes a width look interpretable as a damping
coefficient when it is not, which is why the procedure tests the magnitude and
not the sign.

\section{Coherence loss and phase mixing beyond this scope}

Two further observations bear on Step~2 but are stated at the level of physics rather than as worked entries, because the present package does not establish them at the standard the engine claims require.

First, the reversibility of coherence loss is not binary. Static inhomogeneity is perfectly refocused by an echo; pure dephasing from rapid stochastic frequency fluctuation is largely not refocused; and intermediate correlation times are partly refocused. An echo therefore diagnoses the correlation time of the frequency fluctuation, and the clean statement that ensemble desynchronization is reversible holds only in the static limit. Step~2 is defined by the lifetime-versus-dephasing ratio $R_\phi$, not by reversibility, and the echo enters as a correlation-time diagnostic rather than as the definition of a coordinate.

Second, collisionless phase mixing in a plasma (Landau damping) is an instance of reversible phase mixing in the same sense as static inhomogeneity: the phase information is preserved in the particle velocity distribution and can be recovered, as the plasma-echo experiments demonstrate \citep{gould1967,malmberg1968}. An idealized free-streaming calculation reproduces the reversal, but it is not a self-consistent Vlasov--Poisson solution and does not reproduce the Landau pole; it is an analogy that illustrates why a decaying observable need not correspond to irreversible energy loss, not a computational reassignment of Landau damping. Plasma is discussed here as an analogy and is not entered as a worked entry.

\section{The Barrier-Height/Rate Atlas}

The scope boundary of Section~3 identifies what local mode architecture
cannot supply: a reaction rate depends on a global barrier, not only on the
curvature at a well. That boundary is a refusal, not an endpoint. This
section reports a second, structurally separate atlas built to supply the
missing piece honestly: an activation barrier and a transmission factor,
each independently sourced, checked against an independently measured rate.

\subsection{The admission rule}

A candidate reaction enters this atlas only if every one of the following is
stated, not inferred:

\begin{itemize}
\item an explicit barrier value and its exact type (\(\Delta G^{\ddagger}\),
\(\Delta H^{\ddagger}\), \(\Delta E^{\ddagger}\), or an Arrhenius \(E_a\)),
with its standard state and reference temperature;
\item the barrier's provenance, and specifically whether it is independent
of the rate it will be compared against. A barrier back-calculated from the
same measured rate, assuming unit transmission, is an apparent barrier, not
independent evidence: writing \(D=-\ln\kappa_{\rm dyn}\), such a
back-calculated value is \(B_{\rm app}=B+D\), not \(B\), and using it to
``predict'' the rate it was derived from is circular. This atlas refuses
that route on sight;
\item a transmission-factor route appropriate to the barrier type (an Eyring
\(k_BT/h\) prefactor for a \(\Delta G^{\ddagger}\) barrier; a matching
harmonic-transition-state-theory prefactor for a \(\Delta E^{\ddagger}\)
barrier). Combining a free-energy barrier with an additional well-frequency
prefactor double-counts entropy already contained in
\(\Delta G^{\ddagger}\), and is refused;
\item an independently measured or independently computed rate, not
extracted from the same study that supplied the barrier under a shared
undisclosed assumption.
\end{itemize}

Friction at a barrier top is not assumed equal to the damping measured in a
reactant well. The \(\chi\) atlas of Section~4 cannot be substituted directly
into a transmission-factor calculation; a barrier-local generator, where the
regime supports one, is a separate, separately sourced quantity. As of this
release, no admitted coordinate uses a \(\chi\)-atlas value in its rate
calculation: the two atlases are evidentiarily independent, not yet
factually joined on a common reaction, and this document does not claim
otherwise.

\subsection{Coverage}

Table~\ref{tab:barrieratlas} lists one representative coordinate per
admitted reaction family: 26 families, 61 physical coordinates in total
(full detail in the Supplement), spanning all 18 of the operational
mechanistic classes defined for this atlas and seven physical environments
(gas, aqueous solution, organic solution, enzyme active site, surface, solid
state, and a cross-phase comparison architecture exempted from the
environment count as a comparison method rather than a physical
environment). Evidence grade A is reserved for coordinates with a
prospectively fixed, condition-matched, non-circular route and a pointwise
experimental comparator; grade B covers admitted routes with a stated
residual independence limitation (same-study comparison, a processed rather
than raw comparator, a near-temperature extrapolation); grade C covers
reconstructions, proxies, and series-level comparisons that do not resolve
to a raw pointwise rate. No grade was inflated to fill a coverage target:
several environments (surface, solid state) are carried entirely by grade C
evidence, stated as such throughout.

\begin{table}[htbp]
\centering
\noindent\scriptsize\setlength{\tabcolsep}{2pt}\begin{tabular}{@{}p{0.205\linewidth}p{0.085\linewidth}p{0.125\linewidth}p{0.05\linewidth}p{0.12\linewidth}p{0.135\linewidth}@{}}
\toprule
Reaction & Env. & Mech.\ class & Grade & \(\ln(k_{\rm pred}/k_{\rm obs})\)\(^{\dagger}\) & Source \\
\midrule
4-oxalocrotonate tautomerase stepwise two... & enzyme active site & proton transfer & B & +4.12 & \citep{ba_ruizpernia20094otqmmm} \\
5-hexenyl radical $\to$ cyclopentylmethyl ra... & organic solution & addition & B & +4.79 & \citep{ba_gansauer2013hexenyl,ba_chatgilialoglu1981hexenylexperiment} \\
1,3-butadiene + N-methylmaleimide $\to$ Diel... & aqueous solution & pericyclic multibond reorganization & B & +2.28 & \citep{ba_gui2026lipred,ba_gui2025lipreddataset} \\
{}[Fe(H$_2$O)$_6$](2+) + [Fe(H$_2$O)$_6$](3+) electron self... & aqueous solution & electron transfer & B & +0.42 & \citep{ba_koczorbenda2023et,ba_rossorustad2000feet} \\
Cl(-) + CH$_3$Br $\to$ CH$_3$Cl + Br(-) & aqueous solution & substitution displacement & B/C & +0.20 & \citep{ba_campeggio2020,ba_bathgate1959} \\
CH$_3$S(-) + CH$_3$SSCH$_3$ $\to$ CH$_3$SSCH$_3$ + CH$_3$S(-) (ide... & aqueous solution & substitution displacement & C & n/a & \citep{ba_campeggio2020} \\
beta-cyclodextrin + 1-propanol $\to$ beta-cy... & aqueous solution & association complexation & B & +0.86 & \citep{ba_tangchang2018betacd,ba_betacdhistoricalkineticsaggregate} \\
1-ethoxyisoquinoline $\to$ 1-isoquinolone + ... & gas & elimination & C & +12.85 & \citep{ba_mahmoud2023elimination,ba_alawaditaylor1986elimination} \\
surface formate decomposition on Au(110) ... & surface & surface interfacial elementary reactions & C & +3.63 & \citep{ba_monkhorst2022formateaucu} \\
CH$_3$NC $\to$ CH$_3$CN & gas & rearrangement isomerization & A & +0.17 & \citep{ba_nguyen2018,ba_schneider1962nist} \\
cyclobutane $\to$ ethylene + ethylene & gas & dissociation fragmentation & A & +0.76 & \citep{ba_sirjean2006,ba_genaux1953nist} \\
cis-1,3-butadiene + ethene $\to$ cyclohexene & gas & pericyclic multibond reorganization & B & +0.29 & \citep{ba_mcpdftopesf2025} \\
H + C$_2$H$_6$ $\to$ H$_2$ + C$_2$H$_5$ & gas & hydrogen atom transfer & A & +1.28 & \citep{ba_laude2018,ba_ledevillermaux1978nist} \\
H + CH$_4$ $\to$ H$_2$ + CH$_3$ & gas & hydrogen atom transfer & A & +0.60 & \citep{ba_beyer2016,ba_burgessmanion2021nist} \\
H$_2$ + OH $\to$ H$_2$O + H & gas & hydrogen atom transfer & A & +0.55 & \citep{ba_meisnerkastner2016h2oh,ba_orkin2006h2oh} \\
D + H$_2$ $\to$ HD + H & gas & exchange metathesis & A/C & +0.04 & \citep{ba_hankel2006h3pes,ba_westenbergdehaas1967viamielke} \\
vacancy-mediated Li+ lattice hopping in l... & solid state & solid state topochemical transformations & C & +1.79 & \citep{ba_zhou2023lico2fpmd,ba_sugiyama2009lico2musr} \\
collective Li+ conduction in crystalline ... & solid state & solid state topochemical transformations & C & +2.68 & \citep{ba_guo2022lpsaimd,ba_lpshistoricalexperimentaltable} \\
morphinone reductase hydride transfer fro... & enzyme active site & hydride transfer & B & +0.11 & \citep{ba_delgado2017mrqmmm,ba_delgado2017mrexperiment} \\
OH + HO$_2$ $\to$ H$_2$O + O$_2$ & gas & chain propagation termination & B & +2.30 & \citep{ba_speak2023ohho2,ba_chen2024ohho2} \\
Pd(PCy$_3$)$_2$ + PhBr $\to$ trans-Pd(PCy$_3$)$_2$(Ph)(Br) & organic solution & organometallic bond activation migration & B & +3.00 & \citep{ba_mcmullin2014oxidativeaddition,ba_barrioslanderos2009oxidativeaddition} \\
Porphycene trans <$\to$ trans double proton ... & cross phase gas model solution experiment & proton transfer & B & +0.30 & \citep{ba_litman2019porphycene,ba_ciacka2016porphyceneexperiment} \\
V$_6$O$_{13}$(2-) + iAscH- $\to$ V$_6$O$_{11}$(OH)$_2$(2-) + oxid... & organic solution & coupled proton electron transfer & B & +0.05 & \citep{ba_fertigmatson2023pcet} \\
O* + H* $\to$ OH* at B-type Pt(332) step sites & surface & surface interfacial elementary reactions & C & n/a & \citep{ba_nitz2025pt332} \\
(Q-TTF-Q)- intramolecular electron transfer & organic solution & electron transfer & B & +0.22 & \citep{ba_koczorbenda2023et,ba_gautier2003qttfexperiment} \\
SO$_2$(+) + CO $\to$ SO(+) + CO$_2$ & gas & atom or group transfer & B & +0.02 & \citep{ba_catone2019oat,ba_fehsenfeldferguson1973oat} \\
\bottomrule
\end{tabular}
\caption{The twenty-six admitted reaction families of the Barrier-Height/Rate Atlas, one representative coordinate per family (the admitted coordinate with the smallest \(|\ln(k_{\rm pred}/k_{\rm obs})|\); full 61-coordinate detail, including every temperature and replicate, is in the Supplement). Ionic charges are shown in parentheses, e.g.\ (2+), for table-cell compactness. \(^{\dagger}\)A residual of 0 is exact agreement; the sign shows the direction of the discrepancy. Two families (identity thiolate-disulfide exchange, Pt(332) hydrogen oxidation) report an exact quantum benchmark or temperature-series ratio rather than a pointwise residual and are marked n/a here; their own metric is in the Supplement. Grade C rows are reconstructions or proxies, not raw pointwise measurements, and are labelled as such throughout; large residuals are retained rather than discarded, matching the source-gate discipline of Part I.}
\label{tab:barrieratlas}
\end{table}

\subsection{One clean closure, one retained discrepancy}

The isomerization of methyl isocyanide to acetonitrile, CH\(_3\)NC \(\to\)
CH\(_3\)CN, is the atlas's cleanest case and its historical anchor: the
reaction Rabinovitch's group used to help establish RRKM theory itself. An
activation energy computed entirely from first principles, with no
reference to any measured rate (38.25 \(\pm\) 0.25 kcal~mol\(^{-1}\),
coupled-cluster \citep{ba_nguyen2018}), agrees with the classical
high-pressure-limit Arrhenius activation energy from fall-off kinetics
(38.4 kcal~mol\(^{-1}\) \citep{ba_schneider1962nist}) to 0.15
kcal~mol\(^{-1}\), and the resulting thermal rate constants agree with
experiment to within 10\% at every reported temperature. The transmission
factor here is 1 by construction: this is the gas-phase high-pressure
limit, with no condensed-phase friction to model.

Not every admitted case closes this cleanly, and the atlas keeps the ones
that do not. The oxidative addition of bromobenzene to Pd(PCy\(_3\))\(_2\),
using a condition-matched, dispersion-corrected computed free energy
barrier \citep{ba_mcmullin2014oxidativeaddition} against an independently measured apparent
rate \citep{ba_barrioslanderos2009oxidativeaddition}, underpredicts by a factor of 20.14 under
a unit-transmission Eyring baseline. The source's own account of that gap is
retained rather than resolved: an equilibrium with the unreactive
tris-ligated species Pd(PCy\(_3\))\(_3\) competes with the bisphosphine
pathway measured here, and the apparent second-order rate constant reflects
that speciation rather than a single elementary step in isolation. The
coordinate remains grade B for exactly that reason, and the 20-fold residual
is reported, not adjusted away.

\subsection{Held and refused candidates}

The discipline is as visible in what the atlas declines as in what it
admits. A candidate activation free energy for enzymatic hydride transfer,
back-calculated from the same isotopologue rate the atlas would otherwise
use as a closure test, is refused as circular for the reason given above:
a barrier derived from the rate it would validate is not independent of
it. A PCET candidate pairing predicted and experimental rate
constants of different reaction order and different units is refused
outright, with no numerical ratio recorded. A conductivity proxy for
lithium-thiophosphate ionic transport is retained at grade C and explicitly
barred from promotion, because collective transport is not a molecular rate
constant regardless of how well an Arrhenius fit reproduces it. None of
these refusals were softened to raise a coverage count.

\section{Discussion}

The preceding sections state what each safeguard does and what it refuses. This
section addresses a different question, asked repeatedly of the framing rather
than the mathematics: whether the presentation undersells the engine by
describing it mainly through what it prevents. The honest answer is
calibrated, not uniform: in one place the hedged framing is exactly right and
should not be loosened; in another, a genuinely nontrivial result is available
and is stated here explicitly.

\paragraph{What is a result, not only a guardrail.} The certified-multiplicity
boundary of Section~2.3(b) is not merely a safety mechanism sitting in front
of the classifier; it is itself a nontrivial, adversarially derived result.
Order-2 multiplicity is provably decidable from finite-precision evidence
alone, because two distinct eigenvalues never produce a nonzero rank-one
nullity at their mean; no analogous proof exists for order \(m>2\), and two
successive releases of this engine wrongly believed one did, each caught only
by an external counterexample (a linear-roundoff-coincidence carve-out that a
triangular matrix with genuinely distinct, non-roundoff-separated eigenvalues
defeated). The corrected boundary, that no floating-point evidence ever
certifies an order-\(m>2\) defect regardless of how small its members' spread
is, is verified against zero false certifications across perturbation scales
from \(10^{-2}\) to \(10^{-12}\) and orders 2 through 6, both directly and
under random orthogonal similarity. That is a stated, checkable claim about
what finite-precision multiplicity testing can and cannot establish in
general, not a disclaimer attached to a feature.

\paragraph{What remains a bound, not a theorem.} The constructive
counterexample of Section~3 (identical local reactant wells, activated rates
differing by \(e^{5}\approx148\) through the barrier alone) is exactly that: a
constructive counterexample under one stated model, sufficient to keep the
implementation from inferring rate from local architecture, and reproduced as
an executable regression test for that reason. It is not offered, and should
not be read, as a general limitation theorem for local-architecture-based rate
estimators; no such generality is proven. The distinction matters: a
counterexample bounds one specific inference for one stated construction, a
theorem would bound every construction of a stated class. Overstating the
former as the latter is the same error this framework spends Section~2.3 and
Section~4 refusing to make about degeneracy and about a simulated versus
measured anchor, and it is refused here for the same reason.

\paragraph{Reading the linewidth procedure and the worked examples.} The
three-step procedure of Part~III restates a standard decomposition; what it
adds is that each step is executable rather than declarative, refusing with a
stated reason when a step's precondition is not met (Section~5). The Rh(CO)\(_2\)acac
trajectory and the admission table of Section~6 demonstrate that procedure
running end to end on real systems, distinguishing what a source anchors from
what is illustrative at every step; the quantitative claim stays scoped to
the anchored endpoints, as the limitations that follow state plainly. Reframing either as a
general-purpose reproducible pipeline, independent of the specific systems
carried through it, would claim more than either demonstration establishes.

\section{Limitations, investigated}

The single-mode limitations of the engine stand unchanged: a \(\chi\) read from a lone real pole with no partner is undefined (\(\omega_0\) is not recoverable from one pole) and is reported as such rather than invented; mode parentage must be read together with the frequency gap, since a mid-range parentage with an open gap indicates hybridization while a closing gap indicates coalescence, and parentage near one-half does not by itself certify an exceptional point; non-reciprocal (asymmetric \(C\) or \(K\)) systems are rejected at the input contract rather than returned as spurious mechanical modes; and negative damping, non-Hermitian mass, and many-mode (more than three) coupling are untested, with no robustness claimed.

The extension of Part~III adds its own bounded items, stated rather than hidden. Step~1's resolvability threshold is line-shape- and criterion-specific: the equal-Gaussian merger near $\Delta \approx 2\sigma$ is illustrative, not a universal constant, and two unresolved lines are unresolved under a stated line shape and criterion rather than a mathematical continuum. Step~2's $R_\phi = 1$ is the equal-contribution crossover, not a critical point and not an identifiability transition, and is claimed as such: with $T_1$ measured, $\chi$ is recoverable at every $R_\phi$, and $R_\phi$ reports the bias ($1 + R_\phi + R_{\mathrm{or}}$) of the naive width reading rather than a boundary of recoverability. The reversibility of coherence loss is not binary: an echo refocuses static inhomogeneity completely, fast stochastic dephasing largely not, and intermediate correlation times partly, so reversibility diagnoses correlation time and does not by itself define an axis. The single worked example (Rh(CO)$_2$acac) is anchored to the two endpoints the source reports in text (3.4~K: $T_2 = 95.2$~ps, width 0.11~cm$^{-1}$; 250~K: width 1.5~cm$^{-1}$ (source text), from which $T_2 = 7.08$~ps follows by Eq.~1 and is author-derived rather than stated). The intermediate echo and pump-probe points are by-eye transcriptions of the source's Figure~2 and are shown as an illustrative trajectory only; no crossover temperature or intermediate bias factor is reported as a result. The quantitative claim is limited to the endpoint contrast, a homogeneous-width growth of about 13.6; the fitted-lifetime $\chi$ near $3\times10^{-5}$ at both endpoints is source-informed illustrative rather than source-anchored, pending calibrated pixel extraction of the pump-probe points, and the mechanical-placement claim is scoped to this single example with no coupled-pair catalogue claimed. A calibrated redigitization of Figure~2 would be required before the lifetime-derived $\chi$ or bias could be reported quantitatively. The plasma discussion is an analogy, not a worked entry: the free-streaming calculation illustrates reversible phase mixing but is not a self-consistent Vlasov--Poisson solution and does not reproduce the Landau pole.

A separate boundary applies when the generator is not supplied but inferred. Subspace identification from a vector response record recovers the spectrum accurately but does not by itself certify the parent's multiplicity structure: generic finite-data perturbation splits a defective pole, so a fitted generator is a reduced representation in the sense of Section~2.3 and its structural verdict is scoped to the fit. No such estimator is shipped here.

For general first-order nonnormal generators, pole classification does not determine finite-time transient amplification or pseudospectral sensitivity. ChemSA presently reports spectral structure and does not claim a nonmodal stability bound; a pseudospectral or propagator-norm extension remains future work.

\section{Conclusion}

The contribution of this work is an audited classification engine and, built on the trust that engine earns, a hierarchy that determines when a measured linewidth may legitimately be read as a mechanical damping ratio. The engine's three core scientific safeguards (exceptional-point-versus-semisimple discrimination by eigenvector deficiency, provenance gating that scopes an exceptional-point verdict to the supplied representation, and a constructive counterexample demonstrating that local spectral architecture alone cannot supply a reaction rate), together with the three implementation validity gates that keep them trustworthy in practice (scalar-\(\chi\) applicability, source gating of numerical inputs, and finite-precision equality caution), answer questions a published damped-oscillator fit cannot. Around that engine, the damping ratio \(\chi = \gamma/2\omega_0\) supplies a coordinate that reads the same oscillatory-to-monotone boundary wherever a second-order or decoupled modal representation is established, with a qualifying rule that admits only amplitude damping to that coordinate.

Part~III sharpens that qualifying rule into a decision procedure. A measured spectral line is not self-describing: it may combine static inhomogeneous broadening with a homogeneous width, and the homogeneous width may itself contain lifetime relaxation, pure dephasing and an orientational contribution, of which only the lifetime part is a mechanical damping ratio. A three-step procedure separates them: homogeneous isolation; then separation of the lifetime, pure-dephasing, and orientational contributions, requiring an independent lifetime measurement and an orientational term that is measured, bounded, or independently excluded; then mechanical placement, so that a \(\chi\) is computed only from a lifetime certified as such. The ratio \(R_\phi = (1/T_\phi)/(1/2T_1)\) quantifies the bias (\(1 + R_\phi + R_{\mathrm{or}}\)) incurred by reading the total homogeneous width as lifetime damping, with \(R_\phi = 1\) the equal-contribution crossover rather than a boundary of recoverability. The procedure is carried through one worked real system, the CO stretch of Rh(CO)$_2$acac, anchored to the two endpoints the source reports: the homogeneous width grows by a factor of about 13.6 between them (0.11 to 1.5~cm$^{-1}$), and the source establishes that pure dephasing becomes dominant, since the pump-probe lifetime varies far more weakly than the echo coherence time. A source-informed linear reading of the plotted lifetime data places the mechanical damping ratio near $3\times10^{-5}$ at both endpoints, but that numerical value is illustrative until the published lifetime points are calibrated and digitized; the quantitative claim is limited to the width growth, so coherence loss and mechanical damping are demonstrably different in kind even where the damping value itself remains illustrative. The parallel with the engine is exact, in that a degeneracy cannot be classified from eigenvalues alone and a linewidth cannot be classified from its total width alone, and in both cases the honest object is a test rather than a table. The parallel is also computational and not merely thematic: the fitted lifetime at each temperature builds a scalar generator that the Part~I classifier reads back as a stable oscillatory simple mode whose damping coordinate matches the directly computed $\chi$, so the linewidth procedure and the degeneracy engine are demonstrably one pipeline.

\section*{Data and code availability}

The numerical engine requires NumPy and SciPy; figure generation uses Matplotlib and the regression suites use pytest, with tested versions recorded in the release manifest and printed by the package validator. The code runs on a laptop and ships with a reproducibility guide and adversarial test suites. A master validator (\texttt{run\_all\_chemsa.py}) is included; the validator treats a missing file as a failure and requires an explicit success signal from every stage, and reports twenty of twenty stages passing when \texttt{pytest} is available. A separate package validator (\texttt{validate\_package.py}) checks the release as a whole in ten stages, including that every shipped file except \texttt{MANIFEST.json} itself carries a role and a SHA-256 hash in that manifest (the archive checksum authenticates the manifest), and that no undeclared file (a stray runtime artifact, for instance) is present in the release. The biorthogonal response layer is \texttt{biorthogonal\_response.py} with \texttt{test\_biorthogonal.py}; with the crowded-spectrum certification suite (\texttt{test\_crowded\_spectrum.py}) and the provenance source-gate suite (\texttt{test\_provenance\_gate.py}), the engine and safeguard suites comprise 413 tests, all passing. Randomized statistics reproduce from \texttt{test\_randomized.py} with a fixed seed, the external exceptional-point cross-check from \texttt{benchmark\_ep\_pendulum.py}, and the over-damped CsPbBr\(_3\) cell from \texttt{verify\_cspbbr3\_cell.py}. The source gate (\texttt{provenance\_gate.py}) makes an un-sourced number un-usable and is exercised on the worked example: the reported endpoint width growth is computed through the gate from the two source-anchored widths, so an illustrative point cannot enter it. The single worked example of Part~III and its figure reproduce from \texttt{linewidth\_hierarchy.py}, which reads points digitized from the cited source's Figure~2 (shipped as \texttt{rhco2acac\_fig2\_raw\_digitized.csv}), fits the source-described linear $2T_1(T)$ relation, and derives the decomposition and $\chi$. Every internally generated numerical result and every reported source-derived transformation reproduces from the provided scripts and declared input files; none is asserted from memory. Several application scripts begin from literature inputs transcribed by hand into the shipped data files, and those transcriptions are the declared inputs rather than script-generated quantities.


% System-3/H-Ru material is tracked separately until its evidence gate permits
% integration into the main narrative. See notes/SYSTEM3_H_RU0001_MANUSCRIPT_MATERIAL.md.

\bibliographystyle{unsrtnat}
\begin{thebibliography}{99}

\bibitem[Even et al.(2024)]{even2024} Even, N., Nennig, B., Lefebvre, G., Perrey-Debain, E. Experimental observation of exceptional points in coupled pendulums. \textit{J. Sound Vib.} \textbf{575}, 118239 (2024).

\bibitem[Caughey and O'Kelly(1965)]{caughey} Caughey, T. K., O'Kelly, M. E. J. Classical normal modes in damped linear dynamic systems. \textit{J. Appl. Mech.} \textbf{32}, 583 (1965).

\bibitem[Rector and Fayer(1998)]{rectorfayer} Rector, K. D., Fayer, M. D. Vibrational dephasing mechanisms in liquids and glasses: vibrational echo experiments. \textit{J. Chem. Phys.} \textbf{108}, 1794 (1998).
\bibitem[Fransson et al.(2023)]{fransson} Fransson, E., Rosander, P., Eriksson, F., Rahm, J. M., Tadano, T., Erhart, P. Limits of the phonon quasi-particle picture at the cubic-to-tetragonal phase transition in halide perovskites. \textit{Commun. Phys.} \textbf{6}, 173 (2023).

\bibitem[van der Post et al.(2015)]{water2015} van der Post, S. T., Hsieh, C.-S., Okuno, M., Nagata, Y., Bakker, H. J., Bonn, M., Hunger, J. Strong frequency dependence of vibrational relaxation in bulk and surface water reveals sub-picosecond structural heterogeneity. \textit{Nat. Commun.} \textbf{6}, 8384 (2015).
\bibitem[Cannemeijer and de Vries(1973)]{cannemeijer1973} Cannemeijer, F., de Vries, A. E. Relaxation of the bending vibration of CO\(_2\) in pure CO\(_2\) and in mixtures of CO\(_2\) with noble gases. \textit{Physica} \textbf{64}, 123--133 (1973). DOI 10.1016/0031-8914(73)90118-3. The same authors published a companion study of the asymmetric stretch in the same journal and year (\textit{Physica} \textbf{70}, 135, 1973) whose relaxation times are different throughout; the values used here are those of the bending study.
\bibitem[Cannemeijer and de Vries(1974)]{cannemeijer1974} Cannemeijer, F., de Vries, A. E. Vibrational relaxation of CO\(_2\) in CO\(_2\)-N\(_2\) and CO\(_2\)-O\(_2\) mixtures. \textit{Physica} \textbf{74}, 175--183 (1974). DOI 10.1016/0031-8914(74)90193-1.
\bibitem[de Vasconcelos and de Vries(1977)]{ch4vt} de Vasconcelos, M. H., de Vries, A. E. Vibrational relaxation time measurements in CH\(_4\) and CH\(_4\)-rare gas mixtures. \textit{Physica A} \textbf{86}, 490--512 (1977). DOI 10.1016/0378-4371(77)90091-7.
\bibitem[Cottrell et al.(1966)]{cottrell1966} Cottrell, T. L., Macfarlane, I. M., Read, A. W., Young, A. H. Measurement of vibrational relaxation times by the spectrophone. Application to CH\(_4\), CO\(_2\), N\(_2\)O, COS, NH\(_3\) and HCN. \textit{Trans. Faraday Soc.} \textbf{62}, 2655--2666 (1966). DOI 10.1039/TF9666202655.
\bibitem[Dubajic et al.(2025)]{dubajic2025} Dubajic, M., Pusch, A., Nielsen, M. P., Stranks, S. D. On the accurate determination of phonon lifetimes from the spectral domain. \textit{PRX Energy} \textbf{4}, 033020 (2025). DOI 10.1103/vklp-9bvg.
\bibitem[Nennig and Perrey-Debain(2020)]{eastereig} Nennig, B., Perrey-Debain, E. A high-order continuation method to locate exceptional points and to compute Puiseux series with applications to acoustic waveguides. \textit{J. Comput. Phys.} \textbf{412}, 109425 (2020). Software: \texttt{EasterEig}.
\bibitem[H\"anggi et al.(1990)]{hanggi1990} H\"anggi, P., Talkner, P., Borkovec, M. Reaction-rate theory: fifty years after Kramers. \textit{Rev. Mod. Phys.} \textbf{62}, 251 (1990).
\bibitem[Nennig and Even(2023)]{evencode} Nennig, B., Even, N. real-valued-ep2: a library to locate an EP2 for real-valued parameters. GPL-3.0. \texttt{https://github.com/nicolase7en/real-valued-ep2} (2023).
\bibitem[Steinel et al.(2004)]{steinel2004} Steinel, T., Asbury, J. B., Zheng, J., Fayer, M. D. Watching hydrogen bonds break: a transient absorption study of water. \textit{J. Phys. Chem. A} \textbf{108}, 10957 (2004).
\bibitem[Owrutsky et al.(1991)]{owrutsky1991} Owrutsky, J. C., Kim, Y. R., Li, M., Sarisky, M. J., Hochstrasser, R. M. Determination of the vibrational energy relaxation time of the azide ion in protic solvents by two-color transient infrared spectroscopy. \textit{Chem. Phys. Lett.} \textbf{184}, 368--374 (1991). DOI 10.1016/0009-2614(91)80002-F.
\bibitem[Mori et al.(2026)]{mori2026} Mori, T., Maczka, M., Kojima, S. Raman inactive phonon-polariton dispersion of quantum paraelectric KTaO\(_3\) proved by broadband terahertz time-domain spectroscopy and FTIR. \textit{Solids} \textbf{7}, 29 (2026). DOI 10.3390/solids7030029.
\bibitem[Posch et al.(2007)]{posch2007} Posch, Th., Baier, A., Mutschke, H., Henning, Th. Carbonates in space: the challenge of low-temperature data. \textit{Astrophys. J.} \textbf{668}, 993--1000 (2007). arXiv:0706.3963.
\bibitem[Tokmakoff et al.(1994)]{tokmakoff1994} Tokmakoff, A., Sauter, B., Fayer, M. D. Temperature-dependent vibrational relaxation in polyatomic liquids: Picosecond infrared pump-probe experiments. \textit{J. Chem. Phys.} \textbf{100}, 9035 (1994).
\bibitem[Tokmakoff and Fayer(1995)]{tokmakoff1995} Tokmakoff, A., Fayer, M. D. Homogeneous vibrational dynamics and inhomogeneous broadening in glass-forming liquids: Infrared photon echo experiments from room temperature to 10 K. \textit{J. Chem. Phys.} \textbf{103}, 2810 (1995).
\bibitem[Christensen(2026)]{stabilityarch} Christensen, N. Exceptional-point stability boundaries from quantum dissipation to cosmological acceleration. \textit{Sci. Rep.} (2026). DOI 10.1038/s41598-026-56887-7.
\bibitem[Tisseur and Meerbergen(2001)]{tisseur} Tisseur, F., Meerbergen, K. The quadratic eigenvalue problem. \textit{SIAM Rev.} \textbf{43}, 235 (2001).

\bibitem[Keldysh(1951)]{keldysh} Keldysh, M. V. On the characteristic values and characteristic functions of certain classes of non-self-adjoint equations. \textit{Dokl. Akad. Nauk SSSR} \textbf{77}, 11 (1951).

\bibitem[Gould et al.(1967)]{gould1967} Gould, R. W., O'Neil, T. M., Malmberg, J. H. Plasma wave echo. \textit{Phys. Rev. Lett.} \textbf{19}, 219 (1967).

\bibitem[Malmberg et al.(1968)]{malmberg1968} Malmberg, J. H., Wharton, C. B., Gould, R. W., O'Neil, T. M. Plasma wave echo experiment. \textit{Phys. Rev. Lett.} \textbf{20}, 95 (1968).

\bibitem[Al-Awadi(1986)]{ba_alawaditaylor1986elimination} Al-Awadi, N.; Taylor, R. The mechanism of thermal eliminations. Part 21. Rate data for pyrolysis of 2-ethoxyquinoline, 1- and 3-ethoxyisoquinoline, and 1-ethoxythiazole. J. Chem. Soc., Perkin Trans. 2 1986, 1589-1592. DOI 10.1039/P29860001589.
\bibitem[Barrios-Landeros and Carrow(2009)]{ba_barrioslanderos2009oxidativeaddition} Barrios-Landeros, F.; Carrow, B. P.; Hartwig, J. F. Effect of ligand steric properties and halide identity on oxidative addition of haloarenes to Pd(PR3)2. J. Am. Chem. Soc. 2009, 131, 8141-8154. DOI 10.1021/ja900798s.
\bibitem[Bathgate(1959)]{ba_bathgate1959} Bathgate, R. H.; Moelwyn-Hughes, E. A. The Kinetics of Certain Ionic Exchange Reactions. J. Chem. Soc. 1959, 2642-2648. DOI 10.1039/JR9590002642.
\bibitem[Historical aggregate, in Tang and Chang(2018)]{ba_betacdhistoricalkineticsaggregate} Historical experimental beta-cyclodextrin association-rate measurements reproduced with citations in Tang and Chang 2018, Table 2.
\bibitem[Beyer et al.(2016)]{ba_beyer2016} Beyer, A. N.; Richardson, J. O.; Knowles, P. J.; Rommel, J.; Althorpe, S. C. Quantum Tunneling Rates of Gas-Phase Reactions from On-the-Fly Instanton Calculations. J. Phys. Chem. Lett. 2016, 7, 4374-4379. DOI 10.1021/acs.jpclett.6b02115.
\bibitem[Burgess(2021)]{ba_burgessmanion2021nist} Burgess, D. R., Jr.; Manion, J. A. A Chemical Kinetic Mechanism for H Atom Reactions with Saturated Hydrocarbons Containing C1-C4 Atoms. J. Phys. Chem. Ref. Data 2021, 50, 023102; NIST kinetics record for H + CH4. DOI 10.1063/5.0045228.
\bibitem[Campeggio et al.(2020)]{ba_campeggio2020} Campeggio, J.; Bortoli, M.; Orian, L.; Zerbetto, M.; Polimeno, A. Multiscale Modeling of Reaction Rates: Application to Archetypal SN2 Nucleophilic Substitutions. Phys. Chem. Chem. Phys. 2020, 22, 3455-3465. DOI 10.1039/C9CP03841H.
\bibitem[Catone et al.(2019)]{ba_catone2019oat} Catone, D.; Satta, M.; Cartoni, A.; et al. Gas Phase Oxidation of Carbon Monoxide by Sulfur Dioxide Radical Cation: Reaction Dynamics and Kinetic Trend With the Temperature. Front. Chem. 2019, 7, 140. DOI 10.3389/fchem.2019.00140.
\bibitem[Chatgilialoglu and Ingold(1981)]{ba_chatgilialoglu1981hexenylexperiment} Chatgilialoglu, C.; Ingold, K. U.; Scaiano, J. C. Rate constants and Arrhenius parameters for the 5-exo cyclization of the 5-hexenyl radical. J. Am. Chem. Soc. 1981. DOI 10.1021/ja00416a008.
\bibitem[Chen et al.(2024)]{ba_chen2024ohho2} Chen, I.-Y.; Chang, C.-W.; Fittschen, C.; Luo, P.-L. Accurate Kinetic Studies of OH + HO2 Radical-Radical Reaction through Direct Measurement of Precursor and Radical Concentrations with High-Resolution Time-Resolved Dual-Comb Spectroscopy. J. Phys. Chem. Lett. 2024, 15, 3733-3739. DOI 10.1021/acs.jpclett.4c00494.
\bibitem[Ciacka et al.(2016)]{ba_ciacka2016porphyceneexperiment} Ciacka, P.; Fita, P.; Listkowski, A.; Radzewicz, C.; Waluk, J. Evidence for Dominant Role of Tunneling in Condensed Phases and at High Temperatures: Double Hydrogen Transfer in Porphycenes. J. Phys. Chem. Lett. 2016. DOI 10.1021/acs.jpclett.5b02482.
\bibitem[Delgado(2017a)]{ba_delgado2017mrexperiment} Delgado, M.; et al. Experimental limiting hydride-transfer rates for four NADH isotopologues reported alongside independent QM/MM calculations. ACS Catal. 2017, 7, 3190-3198. DOI 10.1021/acscatal.7b00201.
\bibitem[Delgado(2017b)]{ba_delgado2017mrqmmm} Delgado, M.; et al. Convergence of Theory and Experiment on the Role of Preorganization, Quantum Tunneling, and Enzyme Motions into Flavoenzyme-Catalyzed Hydride Transfer. ACS Catal. 2017, 7, 3190-3198. DOI 10.1021/acscatal.7b00201.
\bibitem[Fehsenfeld(1973)]{ba_fehsenfeldferguson1973oat} Fehsenfeld, F. C.; Ferguson, E. E. Gas-phase ion-molecule kinetics measurement for SO2+ + CO, as cited by Catone et al. 2019.
\bibitem[Fertig(2023)]{ba_fertigmatson2023pcet} Fertig, A. A.; Matson, E. M. Connecting Thermodynamics and Kinetics of Proton Coupled Electron Transfer at Polyoxovanadate Surfaces Using the Marcus Cross Relation. Inorg. Chem. 2023, 62, 1958-1967. DOI 10.1021/acs.inorgchem.2c02541.
\bibitem[Gansaeuer et al.(2013)]{ba_gansauer2013hexenyl} Gansaeuer, A.; Seddiqzai, M.; Dahmen, T.; Sure, R.; Grimme, S. Computational study of the rate constants and free energies of intramolecular radical addition to substituted anilines. Beilstein J. Org. Chem. 2013, 9, 1620-1629. DOI 10.3762/bjoc.9.185.
\bibitem[Gautier et al.(2003)]{ba_gautier2003qttfexperiment} Gautier, N.; Dumur, F.; Lloveras, V.; Vidal-Gancedo, J.; Veciana, J.; Rovira, C.; Hudhomme, P. Intramolecular Electron Transfer Mediated by a Tetrathiafulvalene Bridge in a Purely Organic Mixed-Valence System. Angew. Chem. Int. Ed. 2003, 42, 2765.
\bibitem[Genaux and Kern(1953)]{ba_genaux1953nist} Genaux, C. T.; Kern, F.; Walters, W. D. The Thermal Decomposition of Cyclobutane. J. Am. Chem. Soc. 1953, 75, 6196-6199; kinetics record curated by NIST.
\bibitem[Gui(2025)]{ba_gui2025lipreddataset} Gui, L. Dataset for 'Solvent Design Assisted by Mechanistic Insights: Methods and Application to Peptide Synthesis', version 3. Zenodo, 2025; updated 2026. LiPRED-2026 chapter 5 workbook. DOI 10.5281/zenodo.16358189.
\bibitem[Gui et al.(2026)]{ba_gui2026lipred} Gui, L.; Armstrong, A.; Adjiman, C. S.; Sayyed, F. B.; Galindo, A. Generation and Benchmarking of a Diverse Reaction Database of Quantum Mechanical Liquid-Phase Activation Gibbs Free Energies. Phys. Chem. Chem. Phys. 2026, 28, 13007-13020. DOI 10.1039/D6CP00088F.
\bibitem[Guo(2022)]{ba_guo2022lpsaimd} Guo, H.; et al. Artificial Intelligence-Aided Mapping of the Structure-Composition-Conductivity Relationships of Glass-Ceramic Lithium Thiophosphate Electrolytes. Chem. Mater. 2022, 34. DOI 10.1021/acs.chemmater.2c00267.
\bibitem[Hankel et al.(2006)]{ba_hankel2006h3pes} Hankel, M.; Smith, S. C.; Allan, R. J.; Gray, S. K.; Balint-Kurti, G. G. State-to-state reactive differential cross sections for H + H2 -> H2 + H on five potential-energy surfaces. J. Chem. Phys. 2006, 125, 164303. Author-uploaded full text reports 0.417 eV barriers for BKMP2 and CCI. DOI 10.1063/1.2358350.
\bibitem[Koczor-Benda and Mateeva(2023)]{ba_koczorbenda2023et} Koczor-Benda, Z.; Mateeva, T.; Rosta, E. Direct Calculation of Electron Transfer Rates with the Binless Dynamic Histogram Analysis Method. J. Phys. Chem. Lett. 2023, 14, 9935-9942. DOI 10.1021/acs.jpclett.3c02624.
\bibitem[Laude et al.(2018)]{ba_laude2018} Laude, G.; Richardson, J. O.; Rommel, J.; Althorpe, S. C. Instanton Rate Theory Beyond the Harmonic Approximation. Faraday Discuss. 2018, 212, 237-258. DOI 10.1039/C8FD00085A.
\bibitem[Lede(1978)]{ba_ledevillermaux1978nist} Lede, J.; Villermaux, J. Mesure de la constante de vitesse de reaction des atomes d'hydrogene sur l'ethane et le propane en reacteurs tubulaire et parfaitement agite ouverts. Can. J. Chem. 1978, 56; official NIST kinetics record.
\bibitem[Litman et al.(2019)]{ba_litman2019porphycene} Litman, Y.; Richardson, J. O.; Kumagai, T.; Rossi, M. Elucidating the Nuclear Quantum Dynamics of Intramolecular Double Hydrogen Transfer in Porphycene. J. Am. Chem. Soc. 2019, 141, 2526-2534. DOI 10.1021/jacs.8b12471.
\bibitem[Historical aggregate, in Guo et al.(2022)]{ba_lpshistoricalexperimentaltable} Primary experimental lithium-thiophosphate conductivity and activation-energy measurements reproduced with citations in Guo et al. 2022, Table 1.
\bibitem[Mahmoud and Abdel-Rahman(2023)]{ba_mahmoud2023elimination} Mahmoud, M. A. M.; Abdel-Rahman, M. A.; Shibl, M. F. Pyrolytic elimination of ethylene from ethoxyquinolines and ethoxyisoquinolines: a computational study. Sci. Rep. 2023, 13, 6248. DOI 10.1038/s41598-023-33272-2.
\bibitem[McMullin and Fey(2014)]{ba_mcmullin2014oxidativeaddition} McMullin, C. L.; Fey, N.; Harvey, J. N. Computed ligand effects on the oxidative addition of phenyl halides to phosphine supported palladium(0) catalysts. Dalton Trans. 2014, 43, 13545-13556. DOI 10.1039/C4DT01758G.
\bibitem[Seal and Gagliardi(2025)]{ba_mcpdftopesf2025} Seal, A.; Gagliardi, L.; Ferguson, A. L. Computing Reaction Kinetics with MC-PDFT-OPESf: Combining Multireference Electronic Structure Theory and Enhanced Sampling. J. Phys. Chem. Lett. 2025, 16, 11458-11463. DOI 10.1021/acs.jpclett.5c02966.
\bibitem[Meisner(2016)]{ba_meisnerkastner2016h2oh} Meisner, J.; Kastner, J. Reaction Rates and Kinetic Isotope Effects of H2 + OH -> H2O + H. J. Chem. Phys. 2016; open author manuscript arXiv:1605.08776. DOI 10.1063/1.4952649.
\bibitem[Muramoto et al.(2022)]{ba_monkhorst2022formateaucu} Toward benchmarking theoretical computations of elementary rate constants on catalytic surfaces: formate decomposition on Au and Cu. Chem. Sci. 2022, 13. DOI 10.1039/D1SC05127J.
\bibitem[Nguyen et al.(2018)]{ba_nguyen2018} Nguyen, T. L.; Thorpe, J. H.; Bross, D. H.; Ruscic, B.; Stanton, J. F. Unimolecular Reaction of Methyl Isocyanide to Acetonitrile: A High-Level Theoretical Study. J. Phys. Chem. Lett. 2018, 9, 2532-2538. DOI 10.1021/acs.jpclett.8b01259.
\bibitem[Nitz(2025)]{ba_nitz2025pt332} Nitz, F.; et al. An Experimental Benchmark for the Barrier Height of an On-Surface Reaction: Hydrogen Oxidation on Pt(332). J. Am. Chem. Soc. 2025, 147, 37658-37666. DOI 10.1021/jacs.5c12862.
\bibitem[Orkin et al.(2006)]{ba_orkin2006h2oh} Orkin, V. L.; Kozlov, S. N.; Poskrebyshev, G. A.; Kurylo, M. J. Rate Constant for the Reaction of OH with H2 between 200 and 480 K. J. Phys. Chem. A 2006, 110, 6978-6985. DOI 10.1021/jp057035b.
\bibitem[Rosso(2000)]{ba_rossorustad2000feet} Rosso, K. M.; Rustad, J. R. Ab Initio Calculation of Homogeneous Outer Sphere Electron Transfer Rates: Application to M(OH2)6 3+/2+ Redox Couples. J. Phys. Chem. A 2000, 104, 6718-6725.
\bibitem[Ruiz-Pernia et al.(2009)]{ba_ruizpernia20094otqmmm} Ruiz-Pernia, J. J.; Garcia-Viloca, M.; Bhattacharyya, S.; Gao, J.; Truhlar, D. G.; Tunon, I. Critical Role of Substrate Conformational Change in the Proton Transfer Process Catalyzed by 4-Oxalocrotonate Tautomerase. J. Am. Chem. Soc. 2009, 131, 2687-2698. DOI 10.1021/ja8087423.
\bibitem[Schneider(1962)]{ba_schneider1962nist} Schneider, F. W.; Rabinovitch, B. S. The Thermal Unimolecular Isomerization of Methyl Isocyanide. Fall-off Behavior. J. Am. Chem. Soc. 1962, 84, 4215-4230; kinetics record curated by NIST. DOI 10.1021/ja00881a006.
\bibitem[Sirjean et al.(2006)]{ba_sirjean2006} Sirjean, B.; Glaude, P.-A.; Ruiz-Lopez, M. F.; Fournet, R. Detailed Kinetic Study of the Ring Opening of Cycloalkanes by CBS-QB3 Calculations. J. Phys. Chem. A 2006, 110, 12693-12704. DOI 10.1021/jp0651081.
\bibitem[Speak(2023)]{ba_speak2023ohho2} Speak, T. H.; et al. New Measurements and Calculations on the Kinetics of an Old Reaction: OH + HO2 -> H2O + O2. JACS Au 2023, 3, 1684-1694. DOI 10.1021/jacsau.3c00110.
\bibitem[Sugiyama(2009)]{ba_sugiyama2009lico2musr} Sugiyama, J.; et al. A novel tool for detecting Li diffusion in solids containing magnetic ions: mu+SR study on LixCoO2. Phys. Rev. Lett. 2009, 103, 147601. DOI 10.1103/PhysRevLett.103.147601.
\bibitem[Tang(2018a)]{ba_tangchang2018betacd} Tang, Z.; Chang, C.-E. A. Binding Thermodynamics and Kinetics Calculations Using Chemical Host and Guest: A Comprehensive Picture of Molecular Recognition. J. Chem. Theory Comput. 2018, 14, 303-318. DOI 10.1021/acs.jctc.7b00778.
\bibitem[Westenberg(1967)]{ba_westenbergdehaas1967viamielke} Westenberg, A. A.; de Haas, N. J. Chem. Phys. 1967, 47, 1393. Low-temperature D + H2 and H + D2 absolute-rate expressions reproduced in the openly accessible Table I of Mielke et al. (2003).
\bibitem[Zhou(2023)]{ba_zhou2023lico2fpmd} Zhou, Z.; et al. First-Principles Study on the Interplay of Strain and State-of-Charge with Li-Ion Diffusion in the Battery Cathode Material LiCoO2. ACS Appl. Mater. Interfaces 2023, 15, 53614-53622. DOI 10.1021/acsami.3c14444.
\end{thebibliography}


\end{document}
